% ch2.tex -- Module 7, Chapter 2: Light as a wave (diffraction, interference, polarisation)
% Dot points: M7.IQ2.01, M7.IQ2.02, M7.IQ2.03, M7.IQ2.04
% Figures: figures/ch2_*.py (matplotlib, PDF next to each script) and TikZ below.
\chapter{Light as a Wave: Diffraction, Interference and Polarisation}\label{ch2:chap}
\covers{M7.IQ2.01, M7.IQ2.02, M7.IQ2.03, M7.IQ2.04}

\inquiry{What evidence supports the classical wave model of light, and what
predictions can be made using this model?}

\begin{goals}
  \item describe Newton's particle model and Huygens' wave model of light, and
        analyse the evidence for each one
  \item describe and explain the diffraction of light
  \item use $d\sin\theta = m\lambda$ for a double slit and a diffraction grating
  \item use $I = I_\mathrm{max}\cos^2\theta$, and evaluate what polarisation tells us
        about light
\end{goals}

%% ==================================================================
\section{Newton's and Huygens' models of light}\label{ch2:sec:models}
\covers{M7.IQ2.03}

\begin{bigidea}
Around 1700 there were two models of light. Newton said light is a stream of tiny
particles. Huygens said light is a wave. Experiments decided between them.
\end{bigidea}

\begin{keywords}
  \kw{Model}{a simple picture that explains what we see and predicts new results.}
  \kw{Wavefront}{a line joining points of a wave that are in step, such as all the crests.}
  \kw{Normal}{a line at \ang{90} to a surface. Angles of rays are measured from it.}
  \kw{Refractive index ($n$)}{how many times slower light is in a material than in a
    vacuum. No unit.}
\end{keywords}

Nobody can see what light is made of. So scientists build a \term{model}: a simple
picture that explains what we see. If a model's prediction fails, the model must change.

\paragraph{Newton's particle model.} In his book \emph{Opticks} (1704), Isaac Newton said
light is a stream of tiny particles. He called them \term{corpuscles} (``little bodies'').
Think of tiny tennis balls fired out of a lamp. They have mass and obey Newton's laws.
With no force on them they travel in straight lines. This explains sharp shadows.

\begin{sidebyside}[0.62]
\paragraph{Huygens' wave model.} Christiaan Huygens (1678) said light is a wave. He
thought it was a \emph{longitudinal} wave, like sound (this matters in
\cref{ch2:sec:polarisation}). His rule for how a wave moves is \term{Huygens'
principle} (\cref{ch2:fig:huygens}):
\begin{enumerate}
  \item Every point on a \term{wavefront} acts as a tiny new source.
  \item Each point sends out a small circular wave, called a \term{wavelet}.
  \item The line touching the front of all the wavelets is the new wavefront.
\end{enumerate}
\nextside
  \centering
  \begin{tikzpicture}[x=1cm, y=1cm]
    \foreach \y in {-1.0,-0.5,0,0.5,1.0} {
      \draw[hscmuted, line width=0.5pt] (0,\y) ++(-90:0.65) arc (-90:90:0.65);
      \fill[hscink] (0,\y) circle (1.4pt);
    }
    \draw[hscvel, line width=1.1pt] (0,-1.45) -- (0,1.45);
    \node[text=hscvel, anchor=east, font=\footnotesize] at (-0.12,1.25) {wavefront now};
    \draw[hscvel, line width=0.4pt] (-0.12,1.25) -- (-0.02,1.15);
    \draw[hscvel, line width=1.1pt, dashed] (0.65,-1.45) -- (0.65,1.45);
    \node[text=hscvel, anchor=north west, font=\footnotesize] at (0.6,-1.45) {new wavefront};
    \draw[->, hscink, line width=0.7pt] (1.0,0.25) -- (1.6,0.25) node[right, font=\footnotesize] {travel};
    \draw[hscmuted, line width=0.4pt] (1.5,1.05) node[right, font=\footnotesize, text=hscink] {wavelet} -- (0.55,0.85);
  \end{tikzpicture}
  \captionof{figure}{Huygens' principle.}
  \label{ch2:fig:huygens}
\end{sidebyside}

Sound bends around corners, but light seems not to. Huygens could not explain this.
Newton also had a huge reputation. So most scientists followed Newton for about 100
years.

\begin{toolkit}{Sine and parts of a velocity}
In a right-angled triangle, the \emph{hypotenuse} $h$ is the longest side. The side
opposite the angle $\theta$ is $h\sin\theta$. The side next to $\theta$ is $h\cos\theta$.
Draw a velocity $v$ as the hypotenuse, at $\theta$ to the normal. The part along the
normal is next to $\theta$, so it is $v\cos\theta$. The part along the surface is opposite
$\theta$, so it is $v\sin\theta$. To get an angle from its sine, use $\sin^{-1}$:
$\sin\theta = 0.500$ gives $\theta = \ang{30.0}$. Use \textbf{degrees} mode.
\end{toolkit}

\paragraph{Reflection: both models work.} Angles are measured from the \term{normal}, a
line at \ang{90} to the surface. The \term{angle of incidence} ($\theta_\mathrm{i}$) is the
angle between the incoming ray and the normal. The \term{angle of reflection}
($\theta_\mathrm{r}$) is the angle between the reflected ray and the normal
(\cref{ch2:fig:reflection}).
\begin{itemize}
  \item \textbf{Newton:} near the surface a \term{repulsive force} (a push away from the
        surface) acts along the normal. It reverses the part of the corpuscle's velocity
        along the normal. The part along the surface does not change. The corpuscle
        bounces like a perfectly bouncy ball.
  \item \textbf{Huygens:} each point of the surface sends out a wavelet when the
        wavefront reaches it. Points hit earlier have bigger wavelets. The line touching
        all the wavelets is the reflected wavefront. It leaves at the same angle.
\end{itemize}
Both models give $\theta_\mathrm{i} = \theta_\mathrm{r}$. So reflection cannot decide
between them.

\begin{figure}
  \centering
  \begin{tikzpicture}[x=0.8cm, y=0.8cm]
    \begin{scope}
      \fill[hscmuted!22] (-3.4,-0.35) rectangle (3.4,0);
      \draw[hscink, line width=0.8pt] (-3.4,0) -- (3.4,0);
      \draw[guide] (0,0) -- (0,2.15);
      \node[font=\scriptsize, text=hscmuted, anchor=south] at (0,2.15) {normal};
      \coordinate (O) at (0,0); \coordinate (N) at (0,1);
      \coordinate (P1) at (-1.286,1.532); \coordinate (P2) at (1.286,1.532);
      \draw[velocity] (P1) -- (O) node[pos=0.3, left=3pt, text=hscvel] {$v$};
      \draw[velocity] (O) -- (P2) node[pos=0.65, right=3pt, text=hscvel] {$v$};
      \pic[draw=hscink, line width=0.5pt, angle radius=8mm, "$\theta_\mathrm{i}$", angle eccentricity=1.45] {angle = N--O--P1};
      \pic[draw=hscink, line width=0.5pt, angle radius=8mm, "$\theta_\mathrm{r}$", angle eccentricity=1.45] {angle = P2--O--N};
      % velocity parts before (left) and after (right)
      \draw[hscmuted, line width=0.8pt, ->] (-3.1,1.25) -- (-2.52,1.25);
      \draw[hscmuted, line width=0.8pt, ->] (-3.1,1.25) -- (-3.1,0.56);
      \node[font=\scriptsize, anchor=south] at (-2.81,1.3) {before};
      \draw[hscmuted, line width=0.8pt, ->] (2.3,0.25) -- (2.88,0.25) node[right, font=\scriptsize, text=hscink] {same};
      \draw[hscmuted, line width=0.8pt, ->] (2.3,0.25) -- (2.3,0.94) node[above, font=\scriptsize, text=hscink] {reversed};
      \node[font=\scriptsize, anchor=east] at (2.22,0.6) {after};
      \node[font=\footnotesize\bfseries] at (0,-0.85) {(a) Newton: the normal part reverses};
    \end{scope}
    \begin{scope}[xshift=7.4cm]
      \fill[hscmuted!22] (-1.6,-0.35) rectangle (3.4,0);
      \draw[hscink, line width=0.8pt] (-1.6,0) -- (3.4,0);
      \coordinate (A) at (0,0); \coordinate (B) at (1.905,1.1);
      \coordinate (C) at (2.54,0); \coordinate (T) at (0.635,1.1);
      \foreach \cx/\r in {0/1.270, 0.635/0.953, 1.270/0.635, 1.905/0.318}
        \draw[hscmuted, line width=0.5pt] (\cx+\r,0) arc (0:180:\r);
      \draw[hscvel, line width=1.1pt] (A) -- (B);
      \draw[hscvel, line width=1.1pt, dashed] (T) -- (C);
      \draw[ray] (-0.45,0.779) -- (A);
      \draw[ray] (1.455,1.879) -- (B) -- (C);
      \draw[ray] (A) -- ($(A)!1.3!(T)$);
      \draw[ray] (C) -- (2.99,0.779);
      \fill[hscink] (A) circle (1.4pt) node[below=2pt, font=\scriptsize] {A};
      \fill[hscink] (C) circle (1.4pt) node[below=2pt, font=\scriptsize] {C};
      \node[font=\footnotesize\bfseries] at (0.9,-0.85) {(b) Huygens: reflected wavelets};
    \end{scope}
  \end{tikzpicture}
  \caption{Reflection in the two models. (a) Newton: the grey arrows are the parts of the
    velocity along the surface and along the normal. The normal part reverses, so
    $\theta_\mathrm{i} = \theta_\mathrm{r}$. (b) Huygens: the arriving wavefront (solid
    blue) reaches A first and C last. So the wavelet from A is the biggest. The reflected
    wavefront (dashed blue) touches all the wavelets.}
  \label{ch2:fig:reflection}
\end{figure}

\begin{refresher}{Refraction and Snell's law (Year 11)}
\term{Refraction} is when light changes direction as it goes from one material into
another. Light entering water or glass from air bends \emph{towards} the normal. The sheet
gives $n_1\sin\theta_1 = n_2\sin\theta_2$. The values $n_\mathrm{air} = 1.00$ and
$n_\mathrm{water} = 1.33$ are not on the sheet; questions give them. Example: air to water
at \ang{40.0}: $\sin\theta_2 = \dfrac{1.00 \times \sin\ang{40.0}}{1.33} = 0.483$, so
$\theta_2 = \ang{28.9}$.
\end{refresher}

\paragraph{Refraction: the models disagree.} Both models explain why light bends towards
the normal in water. But they need opposite speeds (\cref{ch2:fig:refraction}).
\textbf{Newton:} the water pulls each corpuscle inwards, along the normal. The normal
part of its velocity grows. The part along the surface stays the same. So light should
be \emph{faster} in water. \textbf{Huygens:} the part of the wavefront that enters first
slows down first. So the wavefront swings towards the normal, and light should be
\emph{slower} in water.

\begin{figure}
  \centering
  \begin{tikzpicture}[x=0.8cm, y=0.8cm]
    \begin{scope}
      \fill[hscvel!9] (-2.1,-2.4) rectangle (2.3,0);
      \draw[hscink, line width=0.8pt] (-2.1,0) -- (2.3,0);
      \node[anchor=south west, font=\footnotesize] at (-2.1,0.02) {air};
      \node[anchor=north west, font=\footnotesize] at (-2.1,-0.02) {water};
      \draw[guide] (0,1.5) -- (0,-1.74);
      \node[font=\scriptsize, text=hscmuted, anchor=south] at (0,1.5) {normal};
      \draw[hscmuted, line width=0.6pt, ->] (-1.226,1.028) -- (0,1.028);
      \draw[hscmuted, line width=0.6pt, ->] (0,1.028) -- (0,0.04);
      \draw[velocity] (-1.226,1.028) -- (0,0) node[pos=0.35, left=1pt, text=hscvel] {$v_1$};
      \draw[hscmuted, line width=0.6pt, ->] (0,0) -- (0,-1.740);
      \draw[hscmuted, line width=0.6pt, ->] (0,-1.740) -- (1.226,-1.740);
      \draw[velocity] (0,0) -- (1.226,-1.740) node[pos=0.55, right=1pt, text=hscvel] {$v_2$};
      \node[text=hscmuted, font=\scriptsize, anchor=south east] at (-0.06,1.03) {parallel part};
      \node[text=hscmuted, font=\scriptsize, anchor=north] at (0.62,-1.76) {same parallel part};
      \node[text=hscmuted, font=\scriptsize, anchor=east, align=right] at (-0.05,-1.0) {bigger\\normal part};
      \draw[force] (1.5,0.35) -- (1.5,-0.45) node[right, text=hscforce, font=\footnotesize] {pull};
      \node[font=\footnotesize\bfseries] at (0.1,-3.0) {(a) Newton: faster in water};
    \end{scope}
    \begin{scope}[xshift=6.0cm]
      \fill[hscvel!9] (-2.1,-2.4) rectangle (2.1,0);
      \begin{scope}
        \clip (-2.1,-2.4) rectangle (2.1,1.5);
        \foreach \k in {-7,...,6} {
          \draw[hscvel, line width=0.7pt] (\k*0.718,0) -- ++(50:3.5);
          \draw[hscvel, line width=0.7pt] (\k*0.718,0) -- ++(215.2:3.5);
        }
      \end{scope}
      \draw[hscink, line width=0.8pt] (-2.1,0) -- (2.1,0);
      \draw[guide] (0,1.5) -- (0,-2.4);
      \node[font=\scriptsize, text=hscmuted, anchor=south] at (0,1.5) {normal};
      \draw[ray] (-1.532,1.286) -- (0,0);
      \draw[ray] (0,0) -- (1.094,-1.553);
      \node[fill=white, inner sep=1.5pt, anchor=south west, font=\footnotesize] at (-2.1,1.1) {air};
      \node[fill=hscvel!9, inner sep=1.5pt, anchor=north west, font=\footnotesize] at (-2.1,-1.9) {water};
      \node[font=\footnotesize\bfseries] at (0,-3.0) {(b) Huygens: slower in water};
    \end{scope}
  \end{tikzpicture}
  \caption{Air to water, \ang{50} to \ang{35}. (a) Newton: a pull at the surface makes the
    normal part of the velocity bigger. (b) Huygens: wavefronts (blue) are closer together
    in the slower water, so the ray (orange) turns towards the normal.}
  \label{ch2:fig:refraction}
\end{figure}

In 1850 L\'eon Foucault found that light is \emph{slower} in water
(\cref{ch2:sec:experiments}). This disproved Newton's explanation of refraction and
agreed with Huygens. But on its own it did not prove the wave model.
\Cref{ch2:tab:evidence} collects all the evidence in this chapter.

\begin{table}
  \caption{The evidence for each model of light.}\label{ch2:tab:evidence}
  \centering\small
  \begin{tabular}{@{}L{25mm}L{31mm}L{36mm}L{40mm}@{}}
    \toprule
    \colhead{What we see} & \colhead{Newton (particles)} & \colhead{Huygens (waves)} &
      \colhead{What the evidence shows}\\
    \midrule
    Sharp shadows & particles go straight & hard to explain at first &
      both, once we know light's wavelength is tiny\\
    Reflection & a repulsive force reverses the normal part of the velocity &
      reflected wavelets & both work: no decision\\
    Refraction & faster in water & slower in water & Foucault (1850): slower, so
      Newton is wrong\\
    Diffraction and interference (\cref{ch2:sec:diffraction,ch2:sec:interference}) &
      cannot make dark bands & wavelets spread into the shadow; Young and Fresnel later
      added that waves can cancel & strong support for waves (Young, 1801; Fresnel and
      Arago, 1818)\\
    Polarisation (\cref{ch2:sec:polarisation}) & no natural explanation & Huygens'
      \emph{longitudinal} wave cannot be polarised; a \emph{transverse} wave can & light
      is a transverse wave, so Huygens' model had to change\\
    \bottomrule
  \end{tabular}
\end{table}

\begin{toolkit}{Powers of ten (reminder)}
From \cref{ch1:chap}: $\qty{1}{nm} = \qty{e-9}{m}$, $\qty{1}{\micro m} = \qty{e-6}{m}$,
$\qty{1}{mm} = \qty{e-3}{m}$. So $\qty{550}{nm} = \qty{5.50e-7}{m}$. On a calculator,
use the $\times 10^{x}$ (or EXP) key.
\end{toolkit}

\begin{toolkit}{Rearranging a formula (reminder)}
From \cref{ch1:chap}: to get one symbol on its own, do the same thing to both sides. Example:
$n_\mathrm{x} = \dfrac{c}{v_\mathrm{x}}$. Multiply both sides by $v_\mathrm{x}$:
$n_\mathrm{x}v_\mathrm{x} = c$. Divide both sides by $n_\mathrm{x}$:
$v_\mathrm{x} = \dfrac{c}{n_\mathrm{x}}$. Check by putting numbers back in: both sides
must still be equal.
\end{toolkit}

\begin{formulacard}{Refractive index and the speed of light}
  \cardformula{n_\mathrm{x} = \frac{c}{v_\mathrm{x}}}
  \begin{symbols}
    \sym{$n_\mathrm{x}$}{refractive index of material x}{no unit}
    \sym{$c$}{speed of light in a vacuum, \qty{3.00e8}{m.s^{-1}}}{\si{m.s^{-1}}}
    \sym{$v_\mathrm{x}$}{speed of light in material x}{\si{m.s^{-1}}}
  \end{symbols}
  \inwords{$n$ says how many times slower light is in the material than in a vacuum.}
  \whentouse{you need the speed of light in water or glass, or $n$ from a speed.}
  \whennottouse{to find the angle of a refracted ray: use $n_1\sin\theta_1 = n_2\sin\theta_2$.}
  \rearranged{$v_\mathrm{x} = \dfrac{c}{n_\mathrm{x}}$ \qquad $c = n_\mathrm{x}v_\mathrm{x}$}
  \tensecond{glass, $n = 1.50$: $v = \dfrac{\qty{3.00e8}{m.s^{-1}}}{1.50}
    = \qty{2.00e8}{m.s^{-1}}$.}
  \trap{$n$ is 1 or more, so $v_\mathrm{x}$ is never bigger than $c$. If you get
    $v > c$, you divided the wrong way.}
\end{formulacard}

\begin{workedexample}{Light crossing a swimming pool}\label{ch2:we:pool}
Light travels \qty{25.0}{m} along a pool of water ($n = 1.33$). Find the speed of the
light in the water and the time it takes.
\given{$n = 1.33$, \ $s = \qty{25.0}{m}$, \ $c = \qty{3.00e8}{m.s^{-1}}$ (data sheet)}
\find{the speed $v$ in water, then the time $t$}
\formula{$v_\mathrm{x} = \dfrac{c}{n_\mathrm{x}}$ (rearranged in the toolkit), \quad then
  $t = \dfrac{s}{v}$ (time $=$ distance $\div$ speed; this is $s = ut$ with $a = 0$)}
\substitute{\eqlines{v &= \frac{\qty{3.00e8}{m.s^{-1}}}{1.33} = \qty{2.256e8}{m.s^{-1}}\\
  t &= \frac{\qty{25.0}{m}}{\qty{2.256e8}{m.s^{-1}}} = \qty{1.108e-7}{s}}}
\answer{\final{$v = \qty{2.26e8}{m.s^{-1}}$} and \final{$t = \qty{1.11e-7}{s}$}
  (3 significant figures, like the data)}
\checkstep{$v < c$ \cmark. In a vacuum $t = \qty{8.33e-8}{s}$; the water time is 1.33
  times longer \cmark.}
\end{workedexample}

\begin{workedexample}{Testing the two models with a glass block}\label{ch2:we:glass}
Light in air meets a glass block at \ang{50.0} to the normal. In the glass the ray is at
\ang{30.0} to the normal. (a) Find $n$ for the glass. (b) Find the speed of light in the
glass that the wave model predicts. (c) In Newton's model the part of the velocity along
the surface does not change. Find the speed Newton's model predicts.
\given{$n_1 = 1.00$, \ $\theta_1 = \ang{50.0}$, \ $\theta_2 = \ang{30.0}$, \
  $v_1 = c = \qty{3.00e8}{m.s^{-1}}$}
\find{(a) $n_2$ \quad (b) $v_2$, wave model \quad (c) $v_2$, Newton's model}
\formula{(a) $n_1\sin\theta_1 = n_2\sin\theta_2$ \quad (b) $v_\mathrm{x} = \dfrac{c}{n_\mathrm{x}}$
  \quad (c) Newton (not on the formula sheet): the pull acts only along the normal, so the
  part of the velocity along the surface does not change. From the toolkit, that part is
  $v\sin\theta$, so $v_1\sin\theta_1 = v_2\sin\theta_2$.}
\substitute{\eqlines{&\sin\ang{50.0} = 0.766, \quad \sin\ang{30.0} = 0.500\\
  \text{(a) } n_2 &= \frac{n_1\sin\theta_1}{\sin\theta_2} \quad \text{(divide both sides by $\sin\theta_2$)}\\
  &= \frac{1.00 \times 0.766}{0.500} = 1.532\\
  \text{(b) } v_2 &= \frac{\qty{3.00e8}{m.s^{-1}}}{1.532} = \qty{1.958e8}{m.s^{-1}}\\
  \text{(c) } v_2 &= \frac{v_1\sin\theta_1}{\sin\theta_2} \quad \text{(divide both sides by $\sin\theta_2$)}\\
  &= \frac{\qty{3.00e8}{m.s^{-1}} \times 0.766}{0.500} = \qty{4.596e8}{m.s^{-1}}}}
\answer{(a) \final{$n = 1.53$} \quad (b) \final{$v = \qty{1.96e8}{m.s^{-1}}$} (slower)
  \quad (c) \final{$v = \qty{4.60e8}{m.s^{-1}}$} (faster)}
\checkstep{the models give opposite answers. So measuring the speed of light in glass or
  water decides between them. Foucault's result (slower) matches (b).}
\end{workedexample}

\begin{mistakes}
\begin{itemize}
  \item \xmark\ ``Foucault proved Huygens right.'' \cmark\ It proved Newton's explanation
        of refraction wrong, and agreed with Huygens.
  \item \xmark\ ``Reflection supports the wave model.'' \cmark\ Both models explain it.
  \item \xmark\ ``Newton said light slows down in water.'' \cmark\ His model needs it to
        speed up.
\end{itemize}
\end{mistakes}

\begin{examtip}
``Analyse the evidence for Newton's and Huygens' models'' (5--6 marks). Full marks need:
\begin{enumerate}
  \item each model in one sentence;
  \item for straight-line travel, reflection and refraction, how \emph{each} model
        explains it;
  \item the deciding evidence in time order: interference and diffraction (Young 1801;
        Fresnel and Arago 1818), then Foucault's slower speed in water (1850);
  \item a judgement, such as ``by 1850 the evidence strongly favoured the wave model''.
\end{enumerate}
A table like \cref{ch2:tab:evidence} organises the answer quickly.
\end{examtip}

\begin{quickcheck}
  \item Calculate the speed of light in diamond ($n = 2.42$).
  \item Explain why the law of reflection could not decide between Newton's and
        Huygens' models.
  \item State what each model predicts for the speed of light in water. State what
        Foucault found, and which model's explanation it ruled out.
\end{quickcheck}

%% ==================================================================
\section{Diffraction}\label{ch2:sec:diffraction}
\covers{M7.IQ2.01, M7.IQ2.03}

\begin{bigidea}
Waves spread out as they pass through a gap or around an edge. This is diffraction. The
narrower the gap (compared with the wavelength), the more the wave spreads. Light
diffracts too, which is evidence that light is a wave.
\end{bigidea}

\begin{keywords}
  \kw{Diffraction}{the spreading of a wave as it passes through a gap or around an edge.}
  \kw{Slit width ($a$)}{how wide one slit is. Unit: m.}
  \kw{Central maximum}{the wide bright band in the middle of a single-slit pattern.}
\end{keywords}

\begin{refresher}{Wave words (Year 11)}
A \term{crest} is a highest point of a wave. A \term{trough} is a lowest point. The
\term{amplitude} is the height of a crest above the middle line. The \term{wavelength}
($\lambda$) is the distance from one crest to the next (unit: m). The \term{frequency}
($f$) is the number of waves passing each second (unit: Hz). They are linked by
$v = f\lambda$ (on the sheet); for light in a vacuum $v = c$. Visible light runs from about
\qty{400}{nm} (violet) to \qty{700}{nm} (red). The \term{phase} says where a wave is in
its cycle: at a crest, at a trough, or in between.
\end{refresher}

Stand outside an open door. You can hear people inside, even when you cannot see them.
The sound bends around the edges of the doorway. This is \term{diffraction}: the
spreading of a wave as it passes through a gap or around an edge.

How much a wave spreads depends on the size of the gap compared with the wavelength
(\cref{ch2:fig:gaps}):
\begin{itemize}
  \item \textbf{Gap much wider than $\lambda$:} most of the wave goes straight through.
        Only the edges bend (\cref{ch2:fig:gaps}a).
  \item \textbf{Gap about $\lambda$, or narrower:} the wave spreads out in almost
        circles, far into the ``shadow'' (\cref{ch2:fig:gaps}b).
  \item $\lambda$, $f$ and the speed do not change. Only the direction does.
\end{itemize}

\begin{figure}
  \centering
  \begin{tikzpicture}[x=0.8cm, y=0.8cm]
    \begin{scope}
      \fill[hscink!70] (-0.06,1.3) rectangle (0.06,2.4);
      \fill[hscink!70] (-0.06,-1.3) rectangle (0.06,-2.4);
      \foreach \x in {0.4,0.8,1.2,1.6} {
        \draw[hscvel, line width=0.7pt] (-\x,-2.2) -- (-\x,2.2);
        \draw[hscvel, line width=0.7pt] (\x,-1.3) -- (\x,1.3);
        \draw[hscvel, line width=0.7pt] (0,1.3) ++(0:\x) arc (0:30:\x);
        \draw[hscvel, line width=0.7pt] (0,-1.3) ++(0:\x) arc (0:-30:\x);
      }
      \node[font=\footnotesize\bfseries, align=center] at (0,-3.3) {(a) gap $\gg \lambda$:\\little spreading};
    \end{scope}
    \begin{scope}[xshift=4.3cm]
      \fill[hscink!70] (-0.06,0.2) rectangle (0.06,2.4);
      \fill[hscink!70] (-0.06,-0.2) rectangle (0.06,-2.4);
      \foreach \x in {0.4,0.8,1.2,1.6} {
        \draw[hscvel, line width=0.7pt] (-\x,-2.2) -- (-\x,2.2);
        \draw[hscvel, line width=0.7pt] (0,0) ++(-82:\x) arc (-82:82:\x);
      }
      \node[font=\footnotesize\bfseries, align=center] at (0,-3.3) {(b) gap $\approx \lambda$:\\strong spreading};
    \end{scope}
  \end{tikzpicture}
  \caption{Water waves (seen from above) reaching a gap. Blue lines are crests, one
    wavelength apart. ($\gg$ means ``much bigger than''; $\approx$ means ``about equal
    to''.)}
  \label{ch2:fig:gaps}
\end{figure}

Huygens' principle explains this. Only the points inside the gap send out wavelets. The
wavelets at the edges spread sideways into the shadow.

\paragraph{Why light seems to go in straight lines.} Light's wavelength is less than a
millionth of a metre. A doorway is over a million times wider. So light spreads by an
amount we cannot notice, and shadows look sharp. To see light diffract, you need a slit
narrower than about \qty{0.1}{mm}.

\begin{sidebyside}[0.5]
  \centering
  \includegraphics{ch2_singleslit.pdf}
  \captionof{figure}{Single-slit patterns from a red laser (same peak height).}
  \label{ch2:fig:singleslit}
\nextside
  \paragraph{The single-slit pattern.} A laser through one narrow slit gives a wide,
  bright \term{central maximum}: the bright band in the middle of the pattern. Beside it
  are weak bright bands and dark bands (\cref{ch2:fig:singleslit}). The dark bands form
  where wavelets from different parts of the slit cancel.
  \begin{itemize}
    \item A \textbf{narrower slit} gives a \textbf{wider} pattern.
    \item A \textbf{longer wavelength} gives a \textbf{wider} pattern: red spreads more
          than violet.
  \end{itemize}
  The spreading depends on $\lambda \div a$, where $a$ is the slit width. The HSC course
  needs no single-slit formula. Describe and explain the pattern in words.
\end{sidebyside}

\paragraph{Diffraction and the two models.} Newton suggested that forces at the edges
bend the corpuscles. But particles cannot explain the dark bands. Adding light to a place
cannot make it darker. Waves explain the spreading, the dark bands and why red spreads
more. In 1818 Fresnel's wave theory even predicted a bright spot in the middle of the
shadow of a small disc. Arago looked, and found it.

\begin{workedexample}{Hearing, but not seeing, through a doorway}\label{ch2:we:door}
A doorway is \qty{0.80}{m} wide. Sound of \qty{500}{Hz} travels at \qty{340}{m.s^{-1}}.
Green light has $\lambda = \qty{550}{nm}$. Use $\dfrac{\text{gap}}{\lambda}$ to explain
why the sound spreads into the room beside the door but the light does not.
\given{gap $= \qty{0.80}{m}$; \ sound $f = \qty{500}{Hz}$, $v = \qty{340}{m.s^{-1}}$; \
  light $\lambda = \qty{5.50e-7}{m}$}
\find{$\lambda$ of the sound, then $\dfrac{\text{gap}}{\lambda}$ for each wave}
\formula{$v = f\lambda$, so $\lambda = \dfrac{v}{f}$ (divide both sides by $f$)}
\substitute{\eqlines{\lambda_\mathrm{sound} &= \frac{\qty{340}{m.s^{-1}}}{\qty{500}{Hz}}
  = \qty{0.68}{m}\\
  \text{sound: } \frac{\text{gap}}{\lambda} &= \frac{\qty{0.80}{m}}{\qty{0.68}{m}} = 1.2\\
  \text{light: } \frac{\text{gap}}{\lambda} &= \frac{\qty{0.80}{m}}{\qty{5.50e-7}{m}}
  = \num{1.5e6}}}
\answer{the gap is about \final{1.2} sound wavelengths, so sound diffracts strongly. It is
  about \final{\num{1.5e6}} light wavelengths, so light hardly spreads.}
\checkstep{this matches everyday life: you hear around a door, but the door's shadow looks
  sharp \cmark.}
\end{workedexample}

\begin{workedexample}{Radio behind a hill}\label{ch2:we:radio}
A town lies behind a hill about \qty{500}{m} wide. An AM radio station broadcasts at
\qty{900}{kHz}. An FM station broadcasts at \qty{100}{MHz}. Radio waves travel at
$c = \qty{3.00e8}{m.s^{-1}}$. Explain why the town receives the AM station well but the FM
station badly.
\given{hill width $\approx \qty{500}{m}$; \ AM $f = \qty{9.00e5}{Hz}$; \
  FM $f = \qty{1.00e8}{Hz}$; \ $c = \qty{3.00e8}{m.s^{-1}}$}
\find{$\lambda$ of each wave, then $\dfrac{\text{hill width}}{\lambda}$}
\formula{$v = f\lambda$ with $v = c$, so $\lambda = \dfrac{c}{f}$ (divide both sides by $f$)}
\substitute{\eqlines{\lambda_\mathrm{AM} &= \frac{\qty{3.00e8}{m.s^{-1}}}{\qty{9.00e5}{Hz}}
  = \qty{333}{m}\\
  \lambda_\mathrm{FM} &= \frac{\qty{3.00e8}{m.s^{-1}}}{\qty{1.00e8}{Hz}} = \qty{3.00}{m}\\
  \text{AM: } \frac{\text{hill width}}{\lambda} &= \frac{\qty{500}{m}}{\qty{333}{m}} = 1.5\\
  \text{FM: } \frac{\text{hill width}}{\lambda} &= \frac{\qty{500}{m}}{\qty{3.00}{m}} = 167}}
\answer{the hill is only about \final{1.5} AM wavelengths wide, so AM waves diffract
  strongly around it and reach the town. It is about \final{167} FM wavelengths wide, so
  FM waves hardly bend. The town is in their ``shadow''.}
\checkstep{diffraction is strong when the gap or obstacle is about one wavelength wide, or
  less. AM wavelengths are hundreds of metres, like hills \cmark.}
\end{workedexample}

\begin{mistakes}
\begin{itemize}
  \item \xmark\ ``A wider gap gives more diffraction.'' \cmark\ A narrower gap (nearer to
        $\lambda$) gives more spreading.
  \item \xmark\ ``Diffraction changes the wavelength.'' \cmark\ Only the direction
        changes. $\lambda$, $f$ and the speed stay the same.
  \item \xmark\ ``Light does not diffract, because shadows are sharp.'' \cmark\ It does,
        but $\lambda$ is tiny. You need a slit narrower than about \qty{0.1}{mm} to see it.
\end{itemize}
\end{mistakes}

\begin{examtip}
\begin{itemize}
  \item ``Describe the pattern from a narrow slit'' (2--3 marks): a wide bright central
        band, weaker bands beside it, dark bands between; wider for a narrower slit or a
        longer $\lambda$. Add a labelled sketch.
  \item ``Explain how diffraction supports the wave model'' (3 marks): waves spread
        through gaps (Huygens' wavelets); particles cannot cancel to make dark bands;
        light spreads and makes dark bands, more for a longer $\lambda$, so it behaves
        as a wave.
\end{itemize}
\end{examtip}

\begin{quickcheck}
  \item Water waves meet a gap \qty{2.0}{cm} wide. Which spread out more: waves with
        $\lambda = \qty{1.5}{cm}$ or with $\lambda = \qty{0.30}{cm}$? Explain.
  \item A red laser shines through a narrow slit onto a screen. Describe how the pattern
        changes when (a) the slit is made narrower, (b) the red laser is replaced by a
        blue laser.
  \item Explain why the dark bands in a single-slit pattern are evidence for the wave
        model of light.
\end{quickcheck}

%% ==================================================================
\section{Interference: Young's double-slit experiment}\label{ch2:sec:interference}
\covers{M7.IQ2.02, M7.IQ2.03}

\begin{bigidea}
Where two waves overlap, they add or cancel. This is interference. Light through two
narrow slits makes bright and dark bands on a screen. Only waves can cancel to make
darkness, so Young's experiment is strong evidence that light is a wave.
\end{bigidea}

\begin{keywords}
  \kw{Superposition}{when waves meet, their displacements add.}
  \kw{Constructive / destructive interference}{waves add to a bigger wave / waves cancel.}
  \kw{Fringe}{one bright or dark band in an interference pattern.}
  \kw{Coherent}{waves with the same frequency and a fixed phase difference (the shift
    between their crests never changes).}
  \kw{Monochromatic}{one wavelength (one colour) only.}
  \kw{Path difference ($\delta$)}{how much further one wave travels than the other. Unit: m.}
  \kw{Order ($m$)}{the number of a fringe, counted out from the centre (bright:
    $m = 0, 1, 2, \ldots$; dark: $m = 1, 2, \ldots$).}
\end{keywords}

Where two waves meet, their \term{displacements} (how far each wave is above or below
its middle line) add together. This is \term{superposition}
(\cref{ch2:fig:superposition}).
\begin{itemize}
  \item Two waves are \term{in phase} when their crests line up. Crest meets crest, and
        they make a bigger wave. This is \term{constructive interference}: bright light.
  \item Two waves are \term{\ang{180} out of phase} when one is half a wavelength behind
        the other. Crest meets trough, and the waves cancel. This is
        \term{destructive interference}: darkness.
\end{itemize}

\begin{figure}
  \centering
  \begin{tikzpicture}[x=0.72cm, y=0.8cm]
    \foreach \col/\sgn/\name in {0/1/{(a) constructive}, 3.0/-1/{(b) destructive}} {
      \begin{scope}[xshift=\col cm]
        \draw[hscseriesA, line width=0.9pt, domain=0:3.2, samples=60] plot (\x, {0.25*sin(\x*225)});
        \draw[hscseriesB, line width=0.9pt, domain=0:3.2, samples=60] plot (\x, {-0.8+\sgn*0.25*sin(\x*225)});
        \draw[hscink, line width=1.2pt, domain=0:3.2, samples=60] plot (\x, {-1.85+(1+\sgn)*0.25*sin(\x*225)});
        \node[font=\footnotesize\bfseries] at (1.6,-2.65) {\name};
      \end{scope}
    }
    \node[font=\footnotesize, anchor=east] at (-0.05,0) {wave 1};
    \node[font=\footnotesize, anchor=east] at (-0.05,-0.8) {wave 2};
    \node[font=\footnotesize\bfseries, anchor=east] at (-0.05,-1.85) {sum};
  \end{tikzpicture}
  \caption{Two waves (a) in phase and (b) \ang{180} out of phase.}
  \label{ch2:fig:superposition}
\end{figure}

Light $+$ light $=$ darkness only makes sense for waves. Two streams of particles just
add up to more light.

In about 1801, Thomas Young passed light through two narrow slits that were very close
together (\cref{ch2:fig:young}). On a far screen he saw evenly spaced bright and dark
bands. These bands are called \term{fringes}.

Light diffracts at each slit, so the waves from the two slits overlap. Where they arrive
in phase, there is a bright fringe. Where they arrive \ang{180} out of phase, there is a
dark fringe.

The two slits must be \term{coherent} sources: their waves have the same frequency and a
fixed phase difference. So Young lit both slits from one small opening (sunlight through a
pinhole). Today a single slit or a laser does the same job. Two separate lamps change
phase at random, so their pattern washes out.

\term{Monochromatic} light has one wavelength only. It gives the clearest fringes. A laser
is both monochromatic and coherent.

\begin{figure}
  \centering
  \begin{tikzpicture}[x=0.9cm, y=0.75cm]
    \fill[hsctip!70] (0,0) circle (0.25);
    \draw[hscink, line width=0.6pt] (0,0) circle (0.25);
    \foreach \a in {0,45,...,315} \draw[hsctip, line width=0.7pt] (\a:0.33) -- (\a:0.48);
    \node[font=\footnotesize, align=center, anchor=north] at (0,-0.6) {light\\source};
    \fill[hscink!75] (1.3,0.07) rectangle (1.38,2.3);
    \fill[hscink!75] (1.3,-0.07) rectangle (1.38,-2.3);
    \begin{scope}
      \clip (1.38,-2.2) rectangle (3.0,2.2);
      \foreach \r in {0.291,0.582,...,1.6} \draw[hscmuted!70, line width=0.45pt] (1.34,0) circle (\r);
    \end{scope}
    \node[font=\footnotesize, anchor=south] at (1.34,2.35) {single slit};
    \fill[hscink!75] (3.0,0.67) rectangle (3.08,2.3);
    \fill[hscink!75] (3.0,0.53) rectangle (3.08,-0.53);
    \fill[hscink!75] (3.0,-0.67) rectangle (3.08,-2.3);
    \begin{scope}
      \clip (3.08,-2.3) rectangle (7.0,2.3);
      \foreach \r in {0.291,0.582,...,4.2} {
        \draw[hscmuted!70, line width=0.45pt] (3.04,0.6) circle (\r);
        \draw[hscmuted!70, line width=0.45pt] (3.04,-0.6) circle (\r);
      }
      \foreach \yy in {0,1.0,-1.0,2.0,-2.0} \draw[hscray, line width=0.9pt, dashed] (3.04,0) -- (7.0,\yy);
    \end{scope}
    \node[font=\footnotesize, anchor=south] at (3.04,2.35) {double slit};
    \node[font=\footnotesize, anchor=east] at (2.95,0.85) {$S_1$};
    \node[font=\footnotesize, anchor=east] at (2.95,-0.85) {$S_2$};
    \draw[|<->|, hscink, line width=0.5pt] (2.45,0.6) -- (2.45,-0.6) node[midway, left, font=\footnotesize] {$d$};
    \fill[black!85] (7.0,-2.4) rectangle (7.3,2.4);
    \foreach \yy in {0,1.0,-1.0,2.0,-2.0} \fill[hscforce!80!white] (7.0,\yy-0.17) rectangle (7.3,\yy+0.17);
    \node[font=\footnotesize, anchor=south] at (7.15,2.45) {screen};
    \node[font=\footnotesize, anchor=west] at (7.4,0) {bright ($m = 0$)};
    \node[font=\footnotesize, anchor=west] at (7.4,0.5) {dark};
    \node[font=\footnotesize, anchor=west] at (7.4,1.0) {bright ($m = 1$)};
    \node[font=\footnotesize, anchor=west] at (7.4,1.5) {dark};
    \node[font=\footnotesize, anchor=west] at (7.4,2.0) {bright ($m = 2$)};
    \draw[|<->|, hscink, line width=0.5pt] (3.04,-2.75) -- (7.0,-2.75) node[midway, below, font=\footnotesize] {$L$ (not to scale: really $L \gg d$)};
  \end{tikzpicture}
  \caption{Young's experiment (modern version, with a single slit as the one source).
    Grey circles are crests from each slit. Along the orange dashed lines the waves
    arrive in phase, so bright fringes form. In between, the waves cancel.}
  \label{ch2:fig:young}
\end{figure}

\subsection{Path difference and \texorpdfstring{$d\sin\theta = m\lambda$}{d sin(theta) = m lambda}}
Light from the two slits travels different distances to a point P on the screen. The
difference between these two distances is the \term{path difference} ($\delta$). If
$\delta = 0, \lambda, 2\lambda, \ldots$ the waves arrive in phase: bright. If
$\delta = \tfrac{1}{2}\lambda, 1\tfrac{1}{2}\lambda, \ldots$ they arrive \ang{180} out of
phase: dark. At the centre $\delta = 0$, so the centre is always bright. To find $\delta$
(\cref{ch2:fig:pathdiff}):
\begin{enumerate}
  \item The screen is far away compared with $d$, so the two rays to P are almost parallel.
  \item Draw $S_1$Q at \ang{90} to the rays. From $S_1$ and from Q, the two paths to P
        are equal. So $\delta = S_2$Q.
  \item The rays make angle $\theta$ with the central line. The slit line $S_1S_2$ is at
        \ang{90} to the central line. The line $S_1$Q is at \ang{90} to the rays. So
        $S_1S_2$ and $S_1$Q are the central line and the ray, each turned through
        \ang{90}. Turning both lines by the same amount does not change the angle between
        them. So the angle at $S_1$ is also $\theta$.
  \item In the right-angled triangle $S_1$Q$S_2$ the hypotenuse is $d$, and $\delta$ is
        opposite $\theta$. So $\delta = d\sin\theta$. Bright needs $\delta = m\lambda$.
\end{enumerate}
$\delta = d\sin\theta$ is not on the formula sheet. Learn it, or derive it as above.

\begin{figure}
  \centering
  \begin{tikzpicture}[x=0.85cm, y=0.82cm]
    \begin{scope}
      \fill[hscink!75] (-0.05,0.68) rectangle (0.05,1.3);
      \fill[hscink!75] (-0.05,0.52) rectangle (0.05,-0.52);
      \fill[hscink!75] (-0.05,-0.68) rectangle (0.05,-1.3);
      \draw[hscink, line width=1.2pt] (6,-1.3) -- (6,2.4) node[above, font=\footnotesize] {screen};
      \draw[guide] (0,0) -- (6,0);
      \coordinate (P) at (6,2.0);
      \coordinate (M) at (0,0); \coordinate (X) at (6,0);
      \draw[ray] (0,0.6) -- (P);
      \draw[ray] (0,-0.6) -- (P);
      \draw[hscmuted, line width=0.5pt] (0,0) -- (P);
      \pic[draw=hscink, line width=0.5pt, angle radius=11mm, "$\theta$", angle eccentricity=1.25] {angle = X--M--P};
      \fill[hscink] (P) circle (1.6pt) node[right, font=\footnotesize] {P};
      \draw[|<->|, line width=0.5pt] (6.65,0) -- (6.65,2.0) node[midway, right, font=\footnotesize] {$y$};
      \draw[|<->|, line width=0.5pt] (0,-1.6) -- (6,-1.6) node[midway, below, font=\footnotesize] {$L$};
      \node[font=\footnotesize, anchor=east] at (-0.1,0.75) {$S_1$};
      \node[font=\footnotesize, anchor=east] at (-0.1,-0.75) {$S_2$};
      \node[font=\footnotesize, anchor=north] at (3.2,-0.02) {central line};
      \node[font=\footnotesize\bfseries] at (3,-2.3) {(a) the set-up (not to scale)};
    \end{scope}
    \begin{scope}[xshift=8.4cm, yshift=0.15cm]
      \coordinate (S1) at (0,1); \coordinate (S2) at (0,-1);
      \coordinate (Q) at (0.866,-0.5); \coordinate (H) at (1.2,-1);
      \coordinate (R2) at (3.2,0.848);
      \draw[ray] (S1) -- (2.4,2.386);
      \draw[ray] (S2) -- (R2);
      \draw[hscink, line width=0.8pt] (S1) -- (S2) node[midway, left, font=\footnotesize] {$d$};
      \draw[hscink, dashed, line width=0.7pt] (S1) -- (Q);
      \draw[guide] (S2) -- (2.0,-1);
      \draw[hscforce, line width=2pt] (S2) -- (Q);
      \node[text=hscforce, font=\footnotesize, anchor=north west] at (0.05,-1.12) {$\delta = d\sin\theta$};
      \pic[draw=hscink, line width=0.5pt, angle radius=4mm, "$\theta$", angle eccentricity=1.5] {angle = S2--S1--Q};
      \pic[draw=hscink, line width=0.5pt, angle radius=13mm, "$\theta$", angle eccentricity=1.18] {angle = H--S2--Q};
      \pic[draw=hscink, line width=0.4pt, angle radius=1.6mm] {right angle = S1--Q--R2};
      \fill[hscink] (S1) circle (1.5pt) node[above left, font=\footnotesize] {$S_1$};
      \fill[hscink] (S2) circle (1.5pt) node[below left, font=\footnotesize] {$S_2$};
      \node[font=\footnotesize, anchor=north west, xshift=0.5pt, yshift=0.5pt] at (Q) {Q};
      \node[font=\footnotesize, anchor=west] at (2.45,2.35) {to P};
      \node[font=\footnotesize, anchor=west] at (3.25,0.85) {to P};
      \node[font=\footnotesize\bfseries] at (1.3,-2.45) {(b) close-up of the slits};
    \end{scope}
  \end{tikzpicture}
  \caption{(a) Light from both slits reaches P, a height $y$ up a screen $L$ away.
    (b) Close-up: the extra path for light from $S_2$ is $\delta = d\sin\theta$ (pink).}
  \label{ch2:fig:pathdiff}
\end{figure}

\begin{formulacard}{Bright fringes (double slit or grating)}
  \cardformula{d\sin\theta = m\lambda}
  \begin{symbols}
    \sym{$d$}{slit separation: distance between the centres of neighbouring slits}{\si{m}}
    \sym{$\theta$}{angle from the central line to the bright fringe}{degree (\si{\degree})}
    \sym{$m$}{order of the bright fringe: 0 (centre), 1, 2, \ldots}{no unit}
    \sym{$\lambda$}{wavelength}{\si{m}}
  \end{symbols}
  \inwords{a bright fringe forms where the path difference $d\sin\theta$ is a whole number
    of wavelengths.}
  \whentouse{bright fringes of a double slit or a grating, with a far screen ($L \gg d$).}
  \whennottouse{dark fringes (next card), or with the width of one slit ($a$, not $d$).}
  \rearranged{$\lambda = \dfrac{d\sin\theta}{m}$ \qquad $\sin\theta = \dfrac{m\lambda}{d}$
    \qquad $d = \dfrac{m\lambda}{\sin\theta}$}
  \tensecond{$d = \qty{2.00e-6}{m}$, $\lambda = \qty{500}{nm}$, $m = 1$:
    $\sin\theta = \dfrac{5.00\times10^{-7}}{2.00\times10^{-6}} = 0.250$, $\theta = \ang{14.5}$.}
  \trap{change nm, \si{\micro m} and mm into metres first. Here $m$ is the order, not a
    mass. Calculator in degrees.}
\end{formulacard}

\begin{formulacard}[notonsheet]{Dark fringes}
  \cardformula{d\sin\theta = \left(m - \tfrac{1}{2}\right)\lambda, \qquad m = 1, 2, 3, \ldots}
  \begin{symbols}
    \sym{$m$}{number of the dark fringe, counted from the centre: 1 for the first}{no unit}
    \sym{$d$, $\theta$, $\lambda$}{as in the bright-fringe card}{\si{m}, \si{\degree}, \si{m}}
  \end{symbols}
  \inwords{a dark fringe forms where the path difference is a whole number of wavelengths
    plus a half.}
  \whentouse{the angle of a dark fringe (a minimum) of a \emph{double slit}.}
  \whennottouse{bright fringes, or gratings (a grating is dark almost everywhere between
    its sharp bright lines). Never put $m = 0$ in this formula.}
  \rearranged{$\sin\theta = \dfrac{\left(m - \frac{1}{2}\right)\lambda}{d}$}
  \tensecond{same slits, first dark fringe ($m = 1$): $\sin\theta =
    \dfrac{0.5 \times 5.00\times10^{-7}}{2.00\times10^{-6}} = 0.125$, $\theta = \ang{7.18}$.}
  \derive{dark needs $\delta = \frac{1}{2}\lambda, \frac{3}{2}\lambda, \frac{5}{2}\lambda,
    \ldots = \left(m - \frac{1}{2}\right)\lambda$. Set this equal to $\delta = d\sin\theta$.}
  \trap{some books write $\left(m + \frac{1}{2}\right)\lambda$ with $m = 0, 1, 2, \ldots$
    Count from the centre: the first dark fringe has $\delta = \frac{1}{2}\lambda$ and the
    third has $\delta = 2\frac{1}{2}\lambda$ in both versions.}
\end{formulacard}

%% Added by the Module 7 coverage audit (M7.IQ2.02, path difference from two distances).
\subsection{When you know both distances}

\begin{sidebyside}[0.42]
  \centering
  \begin{tikzpicture}[x=0.9cm, y=0.9cm]
    \fill[hscink!75] (-0.06,0.75) rectangle (0.06,1.6);
    \fill[hscink!75] (-0.06,0.35) rectangle (0.06,-0.35);
    \fill[hscink!75] (-0.06,-0.75) rectangle (0.06,-1.6);
    \coordinate (S1) at (0,0.55); \coordinate (S2) at (0,-0.55);
    \coordinate (P) at (4.3,1.45);
    \draw[ray] (S1) -- (P) node[pos=0.5, above, sloped, font=\footnotesize, text=hscink] {$S_1P$};
    \draw[ray] (S2) -- (P) node[pos=0.55, below, sloped, font=\footnotesize, text=hscink] {$S_2P$ (longer)};
    \fill[hscink] (S1) circle (1.5pt) node[left=2pt, font=\footnotesize] {$S_1$};
    \fill[hscink] (S2) circle (1.5pt) node[left=2pt, font=\footnotesize] {$S_2$};
    \fill[hscink] (P) circle (1.6pt) node[right, font=\footnotesize] {P};
    \draw[|<->|, hscink, line width=0.5pt] (-0.75,0.55) -- (-0.75,-0.55)
      node[midway, left, font=\footnotesize] {$d$};
  \end{tikzpicture}
  \captionof{figure}{Two sources and a point P. The path difference is
    $\delta = S_2P - S_1P$.}\label{ch2:fig:twodist}
\nextside
  Some questions give the distance from each source to a point P. Then you do not need
  an angle. Subtract the two distances to get the path difference $\delta$
  (\cref{ch2:fig:twodist}). Then divide by $\lambda$:
  \begin{itemize}
    \item a whole number (0, 1, 2, \ldots): crest meets crest. Constructive: bright light,
          loud sound, a strong signal;
    \item a whole number plus a half: crest meets trough. Destructive: dark, quiet, a
          weak signal.
  \end{itemize}
  This works for any waves (light, microwaves, sound, water), near the sources or far away.
\end{sidebyside}

\begin{formulacard}[notonsheet]{Path difference from two distances}
  \cardformula{\delta = \left|S_2P - S_1P\right|}
  \begin{symbols}
    \sym{$\delta$}{path difference at the point P}{\si{m}}
    \sym{$S_1P$, $S_2P$}{distance from source 1, and from source 2, to P}{\si{m}}
    \sym{$\lambda$}{wavelength of the waves}{\si{m}}
    \sym{$m$}{a whole number: 0, 1, 2, \ldots\ (constructive); 1, 2, 3, \ldots\ (destructive)}{no unit}
  \end{symbols}
  \inwords{if $\delta$ is a whole number of wavelengths the waves add (constructive); if it is
    a whole number plus a half they cancel (destructive).}
  \whentouse{two coherent sources that send out waves in phase (two slits, two speakers, two
    antennas), when you know both distances to P.}
  \whennottouse{sources that are \ang{180} out of phase: then constructive and destructive swap.
    If you know only an angle, use $d\sin\theta = m\lambda$.}
  \rearranged{constructive: $\delta = m\lambda$, so $\lambda = \dfrac{\delta}{m}$\newline
    destructive: $\delta = \left(m - \tfrac{1}{2}\right)\lambda$, so
    $\lambda = \dfrac{\delta}{m - \frac{1}{2}}$ \ ($m = 1, 2, \ldots$)}
  \tensecond{$S_1P = \qty{2.00}{m}$, $S_2P = \qty{2.75}{m}$, $\lambda = \qty{0.500}{m}$:
    $\delta = \qty{0.75}{m}$ and $\dfrac{\delta}{\lambda} = 1.5$, so destructive.}
  \derive{the two waves leave their sources in step. The wave from the further source travels
    $\delta$ further, so it arrives $\dfrac{\delta}{\lambda}$ cycles behind. A whole number of
    cycles behind: crest still meets crest. Half a cycle more: crest meets trough. For a far
    screen, $\delta = d\sin\theta$ (above).}
  \trap{subtract the two distances to P. Do not use the distance between the sources. Keep
    $\delta$ and $\lambda$ in the same unit (both in cm, or both in m).}
\end{formulacard}

\subsection{Where the fringes land}

\begin{toolkit}{tan, \texorpdfstring{$\tan^{-1}$}{tan-1} and small angles}
In a right-angled triangle, $\tan\theta = \dfrac{\text{opposite}}{\text{adjacent}}$. (The
adjacent side is the side next to $\theta$ that is not the hypotenuse.) In
\cref{ch2:fig:pathdiff}a the side opposite $\theta$ is $y$ and the side next to it is $L$,
so $\tan\theta = \dfrac{y}{L}$. To get the angle back, use $\tan^{-1}$ (SHIFT, then tan):
$\tan\theta = 0.200$ gives $\theta = \ang{11.3}$.

For small angles $\sin\theta \approx \tan\theta$: they differ by 0.4\,\% at \ang{5},
1.5\,\% at \ang{10} and 6.4\,\% at \ang{20}. Below about \ang{10} you may use
$\sin\theta \approx \dfrac{y}{L}$. For bigger angles find
$\theta = \tan^{-1}\!\left(\dfrac{y}{L}\right)$ first, then take $\sin\theta$.
\end{toolkit}

A fringe at angle $\theta$ lands a distance $y$ from the centre of the central bright
fringe. The screen is a distance $L$ from the slits (\cref{ch2:fig:pathdiff}a).

\begin{formulacard}[notonsheet]{Where a fringe lands on the screen}
  \cardformula{y = L\tan\theta}
  \begin{symbols}
    \sym{$y$}{distance on the screen from the centre of the central bright fringe}{\si{m}}
    \sym{$L$}{distance from the slits to the screen (at \ang{90} to the screen)}{\si{m}}
    \sym{$\theta$}{angle of the fringe from the central line}{degree (\si{\degree})}
  \end{symbols}
  \inwords{the bigger the angle, or the further away the screen, the further the fringe is
    from the centre.}
  \whentouse{changing an angle into a position on the screen, or a measured position into
    an angle. It works at any angle.}
  \whennottouse{on its own to find $\lambda$: first find $\theta$, then use
    $d\sin\theta = m\lambda$. For the gap between neighbouring fringes at small angles,
    $\Delta y = \dfrac{\lambda L}{d}$ is quicker.}
  \rearranged{$\theta = \tan^{-1}\!\left(\dfrac{y}{L}\right)$ \qquad $L = \dfrac{y}{\tan\theta}$}
  \tensecond{$L = \qty{2.00}{m}$, $\theta = \ang{5.0}$: $y = 2.00 \times \tan\ang{5.0}
    = \qty{0.175}{m}$.}
  \derive{in the right-angled triangle of \cref{ch2:fig:pathdiff}a,
    $\tan\theta = \dfrac{\text{opposite}}{\text{adjacent}} = \dfrac{y}{L}$. Multiply both
    sides by $L$.}
  \trap{use tan here, not sin. Measure $y$ from the centre of the central bright fringe,
    not from the next fringe.}
\end{formulacard}

\begin{formulacard}[notonsheet]{Fringe spacing (small angles)}
  \cardformula{\Delta y = \frac{\lambda L}{d}}
  \begin{symbols}
    \sym{$\Delta y$}{distance between the centres of neighbouring bright (or dark) fringes}{\si{m}}
    \sym{$y_m$}{distance from the centre to bright fringe number $m$}{\si{m}}
    \sym{$L$}{distance from the slits to the screen}{\si{m}}
    \sym{$\lambda$, $d$}{wavelength, slit separation}{\si{m}}
  \end{symbols}
  \inwords{fringes are further apart for a longer $\lambda$, a further screen, or slits
    closer together.}
  \whentouse{a double slit at small angles (below about \ang{10}).}
  \whennottouse{large angles, as with most gratings: use $d\sin\theta = m\lambda$ and
    $y = L\tan\theta$.}
  \rearranged{$\lambda = \dfrac{d\,\Delta y}{L}$ \qquad $d = \dfrac{\lambda L}{\Delta y}$}
  \tensecond{$\lambda = \qty{600}{nm}$, $L = \qty{2.0}{m}$, $d = \qty{0.40}{mm}$:
    $\Delta y = \dfrac{6.00\times10^{-7} \times 2.0}{4.0\times10^{-4}} = \qty{3.0e-3}{m}$.}
  \derive{\begin{align*}
      d\,\frac{y_m}{L} &= m\lambda && \text{$d\sin\theta = m\lambda$ with $\sin\theta \approx \tfrac{y_m}{L}$}\\
      y_m &= \frac{m\lambda L}{d} && \text{multiply both sides by $\tfrac{L}{d}$}\\
      \Delta y = y_{m+1} - y_m &= \frac{\lambda L}{d} && \text{subtract neighbouring fringes}
    \end{align*}}
  \trap{$d$ is the slit separation, not a slit's width. $\Delta y$ is the gap between
    neighbours, not the distance from the centre.}
\end{formulacard}

\Cref{ch2:fig:pattern} shows the pattern. From the centre outwards: $m = 0$ bright, 1st
dark, $m = 1$ bright, 2nd dark, \ldots. Neighbouring bright fringes are $\Delta y$ apart,
and so are neighbouring dark fringes. Each dark fringe is halfway between two bright
fringes, $\tfrac{1}{2}\Delta y$ from each.

The fringes fade away from the centre. Each slit has its own single-slit pattern
(\cref{ch2:fig:singleslit}), and this pattern dims the outer fringes. In
\cref{ch2:fig:pattern} it is the dashed curve, called the \term{envelope}.

\begin{figure}
  \centering
  \includegraphics{ch2_doubleslit.pdf}
  \caption{A double-slit pattern: the fringes on the screen (top) and their brightness
    (bottom), with each bright fringe's order $m$.}
  \label{ch2:fig:pattern}
\end{figure}

\paragraph{Predictions.} From $\Delta y = \dfrac{\lambda L}{d}$, the fringes spread
further apart for a longer wavelength (red rather than blue), a further screen, or slits
closer together. With white light, the central fringe is white, because $\delta = 0$ for
every colour. Beside it are coloured fringes, with blue on the inside and red on the
outside.

\paragraph{Impact.} Young showed that light can cancel light. Particles cannot do this.
He also measured the wavelengths of colours. But Newton's reputation was huge. The wave
model was widely accepted only after Fresnel and Arago (1818--1819) and Foucault (1850).
(In Module~8, electrons make the same pattern.)

\begin{workedexample}{Fringes from a green laser}\label{ch2:we:green}
Green laser light ($\lambda = \qty{532}{nm}$) passes through two slits \qty{0.250}{mm}
apart onto a screen \qty{3.00}{m} away. Find (a) the angle of the second-order bright
fringe, (b) its distance from the centre, (c) the fringe spacing.
\given{$\lambda = \qty{5.32e-7}{m}$, \ $d = \qty{2.50e-4}{m}$, \ $L = \qty{3.00}{m}$, \ $m = 2$}
\find{(a) $\theta$ \quad (b) $y$ \quad (c) $\Delta y$}
\formula{$d\sin\theta = m\lambda$, so $\sin\theta = \dfrac{m\lambda}{d}$ (divide both
  sides by $d$); \quad $y = L\tan\theta$; \quad $\Delta y = \dfrac{\lambda L}{d}$}
\substitute{\eqlines{\sin\theta &= \frac{2 \times \qty{5.32e-7}{m}}{\qty{2.50e-4}{m}}
  = \num{4.256e-3}, \text{ so } \theta = \ang{0.2439}\\
  y &= \qty{3.00}{m} \times \tan\ang{0.2439} = \qty{1.277e-2}{m} = \qty{12.8}{mm}\\
  \Delta y &= \frac{\qty{5.32e-7}{m} \times \qty{3.00}{m}}{\qty{2.50e-4}{m}}
  = \qty{6.38e-3}{m} = \qty{6.38}{mm}\\
  &\text{(m to mm: multiply by 1000)}}}
\answer{(a) \final{$\theta = \ang{0.244}$} \quad (b) \final{$y = \qty{12.8}{mm}$} \quad
  (c) \final{$\Delta y = \qty{6.38}{mm}$}}
\checkstep{the $m = 2$ fringe is two spacings out: $2 \times 6.38 = \qty{12.8}{mm}$ \cmark.
  The angle is tiny, so the small-angle formula is fine.}
\end{workedexample}

\begin{workedexample}{Where is the third dark fringe?}\label{ch2:we:dark}
Orange light ($\lambda = \qty{610}{nm}$) passes through two slits \qty{0.180}{mm} apart
onto a screen \qty{2.20}{m} away. How far from the centre is the third dark fringe?
\given{$\lambda = \qty{6.10e-7}{m}$, \ $d = \qty{1.80e-4}{m}$, \ $L = \qty{2.20}{m}$, \
  third dark fringe, so $m = 3$}
\find{$y$, the distance of the third dark fringe from the centre}
\formula{$d\sin\theta = \left(m - \tfrac{1}{2}\right)\lambda$, so
  $\sin\theta = \dfrac{\left(m - \frac{1}{2}\right)\lambda}{d}$ (divide both sides by $d$);
  \quad then $y = L\tan\theta$}
\substitute{\eqlines{m - \tfrac{1}{2} &= 3 - 0.5 = 2.5\\
  \sin\theta &= \frac{2.5 \times \qty{6.10e-7}{m}}{\qty{1.80e-4}{m}} = \num{8.472e-3}\\
  \theta &= \sin^{-1}(\num{8.472e-3}) = \ang{0.4854}\\
  y &= \qty{2.20}{m} \times \tan\ang{0.4854} = \qty{1.864e-2}{m}\\
  &= \qty{18.64}{mm} \quad \text{(m to mm: multiply by 1000)}}}
\answer{\final{$y = \qty{18.6}{mm}$} from the centre (3 significant figures)}
\checkstep{$\Delta y = \dfrac{\lambda L}{d} = \qty{7.456}{mm}$. The third dark fringe is
  $2\tfrac{1}{2}$ spacings out: $2.5 \times 7.456 = \qty{18.6}{mm}$ \cmark.}
\end{workedexample}

\begin{workedexample}{Microwaves from two slits}\label{ch2:we:microwave}
Microwaves of wavelength \qty{2.80}{cm} pass through two slits \qty{10.0}{cm} apart in a
metal sheet. A detector is placed at two points. P is \qty{52.0}{cm} from slit 1 and
\qty{59.0}{cm} from slit 2. Q is \qty{48.6}{cm} from slit 1 and \qty{54.2}{cm} from slit 2.
(a) Is the signal strong or weak at P? (b) At Q? (c) The wavelength is changed. Find the
longest wavelength that gives a strong signal at P.
\given{$\lambda = \qty{2.80}{cm}$; \ P: $S_1P = \qty{52.0}{cm}$, $S_2P = \qty{59.0}{cm}$; \
  Q: $S_1Q = \qty{48.6}{cm}$, $S_2Q = \qty{54.2}{cm}$}
\find{(a), (b) the path difference at P and at Q, in wavelengths \quad (c) the longest
  $\lambda$ that makes P constructive}
\formula{$\delta = \left|S_2P - S_1P\right|$, then $\dfrac{\delta}{\lambda}$. Constructive
  needs $\delta = m\lambda$, so $\lambda = \dfrac{\delta}{m}$ (divide both sides by $m$).}
\substitute{\eqlines{\text{(a) } \delta_\mathrm{P} &= \qty{59.0}{cm} - \qty{52.0}{cm} = \qty{7.0}{cm}\\
  \frac{\delta_\mathrm{P}}{\lambda} &= \frac{\qty{7.0}{cm}}{\qty{2.80}{cm}} = 2.5\\
  \text{(b) } \delta_\mathrm{Q} &= \qty{54.2}{cm} - \qty{48.6}{cm} = \qty{5.6}{cm}\\
  \frac{\delta_\mathrm{Q}}{\lambda} &= \frac{\qty{5.6}{cm}}{\qty{2.80}{cm}} = 2.0\\
  \text{(c) } \lambda &= \frac{\delta_\mathrm{P}}{m} = \frac{\qty{7.0}{cm}}{1} = \qty{7.0}{cm}
  \quad\text{(longest $\lambda$: smallest $m$, which is 1)}}}
\answer{(a) $2\tfrac{1}{2}$ wavelengths: crest meets trough, so the signal is \final{weak}
  (destructive: the third minimum, $m = 3$). (b) 2 whole wavelengths: the signal is
  \final{strong} (constructive: the second-order maximum, $m = 2$). (c)
  \final{$\lambda = \qty{7.0}{cm}$}.}
\checkstep{a path difference can never be more than the slit separation. 7.0 and
  \qty{5.6}{cm} are both less than \qty{10.0}{cm} \cmark. Both lengths are in cm, so the units
  cancel in $\dfrac{\delta}{\lambda}$ \cmark.}
\end{workedexample}

\begin{mistakes}
\begin{itemize}
  \item \xmark\ Using a slit's width for $d$. \cmark\ $d$ is the distance between slit centres.
  \item \xmark\ Leaving $\lambda$ in nm or $d$ in mm. \cmark\ Change both to metres first.
  \item \xmark\ Using $m\lambda$ for a dark fringe. \cmark\ Dark: $\left(m - \frac{1}{2}\right)\lambda$,
        $m = 1, 2, \ldots$
  \item \xmark\ ``A bright fringe and the next dark fringe are $\Delta y$ apart.''
        \cmark\ They are $\tfrac{1}{2}\Delta y$ apart.
  \item \xmark\ Using $\sin\theta \approx \dfrac{y}{L}$ at large angles. \cmark\ Above about
        \ang{10}, find $\theta = \tan^{-1}\!\left(\dfrac{y}{L}\right)$, then $\sin\theta$.
\end{itemize}
\end{mistakes}

\begin{examtip}
\begin{itemize}
  \item ``Describe Young's experiment and assess its impact'' (4--5 marks): a labelled
        diagram (source, single slit, double slit $d$, screen $L$); the result (draw the
        brightness graph, \cref{ch2:fig:pattern}); bright where $\delta = m\lambda$, dark where
        $\delta = \left(m - \frac{1}{2}\right)\lambda$; the judgement: strong evidence for
        waves, but not accepted at once because of Newton's reputation.
  \item Calculations (2--3 marks): show the conversion to metres, the rearranged formula,
        the substitution with units, and the answer with a unit.
  \item ``Is P bright or dark (loud or quiet)?'' (2 marks): one mark for
        $\delta = \left|S_2P - S_1P\right|$ divided by $\lambda$; one mark for the conclusion
        with its reason (whole number: constructive; plus a half: destructive).
\end{itemize}
\end{examtip}

\begin{quickcheck}
  \item Explain why the centre of a double-slit pattern is bright for every wavelength,
        and why it is white when white light is used.
  \item Blue light ($\lambda = \qty{450}{nm}$) passes through two slits \qty{0.200}{mm}
        apart onto a screen \qty{2.40}{m} away. Calculate the fringe spacing.
  \item Do the fringes move closer together or further apart? (a) Red light is replaced
        by green. (b) The screen is moved further away. (c) The slits are moved further
        apart.
  \item Light of wavelength \qty{600}{nm} passes through two slits \qty{0.300}{mm} apart
        onto a screen \qty{2.00}{m} away. Find the angle of the second dark fringe and its
        distance from the centre.
  \item Two loudspeakers are connected to one signal generator, so they give out sound in
        phase. The frequency is \qty{850}{Hz} and the speed of sound is \qty{340}{m.s^{-1}}.
        A microphone is \qty{3.00}{m} from one speaker and \qty{3.60}{m} from the other.
        (a) Is the sound loud or quiet at the microphone? Explain. (b) Find the lowest
        frequency that would make the sound loud there.\label{ch2:qc:speakers}
\end{quickcheck}

%% ==================================================================
\section{Diffraction gratings}\label{ch2:sec:gratings}
\covers{M7.IQ2.02}

\begin{bigidea}
A diffraction grating is a slide with many equally spaced slits. It makes sharp, bright
lines at the angles where $d\sin\theta = m\lambda$. Each colour goes to its own angle, so
a grating splits white light into a spectrum.
\end{bigidea}

\begin{keywords}
  \kw{Diffraction grating}{a slide with many equally spaced, parallel slits.}
  \kw{Grating spacing ($d$)}{the distance from one slit to the next. Unit: m.}
  \kw{Spectrum}{a band of all the colours, in order of wavelength.}
\end{keywords}

A CD shows rainbow colours because its tracks are thousands of close, equally spaced
lines. It acts as a \term{diffraction grating}: a slide with many equally spaced slits.
Real gratings have about 100 to 1000 slits per millimetre. The rule for the bright lines
is still $d\sin\theta = m\lambda$. Here $d$ is the \term{grating spacing}: the distance
from one slit to the next.
\begin{itemize}
  \item \textbf{Sharp, bright lines.} A bright line needs the waves from \emph{all} the
        slits to be in step. This happens only at the exact angles where
        $d\sin\theta = m\lambda$. Just off these angles, the many waves cancel. More slits
        also pass more light.
  \item \textbf{Large angles.} $d$ is tiny, so the lines are far apart and easy to measure.
        Use $d\sin\theta = m\lambda$ and $y = L\tan\theta$, not $\Delta y = \dfrac{\lambda L}{d}$.
  \item \textbf{A highest order.} $\sin\theta$ cannot be more than 1, so $m\lambda$ cannot be
        more than $d$. The highest order is $\dfrac{d}{\lambda}$ rounded \emph{down}. If
        $\dfrac{d}{\lambda}$ is exactly a whole number, that order would be at \ang{90}
        (along the grating) and cannot be seen. Then the highest order you can see is one
        less.
\end{itemize}

\begin{formulacard}[notonsheet]{Grating spacing from lines per metre}
  \cardformula{d = \frac{1}{N}}
  \begin{symbols}
    \sym{$d$}{grating spacing (distance from one slit to the next)}{\si{m}}
    \sym{$N$}{number of lines (slits) in each metre of the grating}{\si{m^{-1}}}
  \end{symbols}
  \inwords{$N$ lines in each metre means the lines are $\dfrac{1}{N}$ of a metre apart.}
  \whentouse{a grating is given in ``lines per mm'' or ``lines per cm''.}
  \whennottouse{$d$ is already given as a distance.}
  \rearranged{$N = \dfrac{1}{d}$}
  \tensecond{250 lines per mm $= \num{2.50e5}$ lines per m, so
    $d = \dfrac{1}{2.50\times10^{5}} = \qty{4.00e-6}{m}$.}
  \derive{$N$ gaps of width $d$ fill one metre, so $Nd = \qty{1}{m}$; divide both sides by $N$.}
  \trap{change to lines per metre first: lines per mm $\times 1000$, lines per cm
    $\times 100$. $d$ is not $N$.}
\end{formulacard}

\paragraph{White light through a grating} (\cref{ch2:fig:spectra}). At $m = 0$ every
colour has $\theta = 0$, so the central line is white. Each order ($m = 1, 2, \ldots$) is
a full \term{spectrum}: a band of all the colours. In each band, violet (shortest
$\lambda$) is nearest the centre and red is furthest out. This is true on both sides.
Higher orders spread more and can overlap. (A prism bends violet most; a grating sends red
to the biggest angle.)

\term{Spectroscopes} (instruments that spread light into a spectrum so you can study it)
use gratings to study light from stars and gases (\cref{ch1:sec:spectra}).

\begin{figure}
  \centering
  \includegraphics{ch2_spectra.pdf}
  \caption{White light (\qtyrange{400}{700}{nm}) through a grating with 300 lines per mm.
    Each order is drawn in its own row. In every order violet is nearest the centre and
    red is furthest out. Orders 2 and 3 overlap between about \ang{21} and \ang{25}.}
  \label{ch2:fig:spectra}
\end{figure}

\begin{workedexample}{A green laser and a grating}\label{ch2:we:grating}
Green laser light ($\lambda = \qty{532}{nm}$) shines through a grating with 500 lines per
mm. Find (a) the angle of the first-order line, (b) the highest order, (c) the number of
bright spots on the wall.
\given{$N = 500 \times 1000 = \num{5.00e5}$ lines per m, \ $\lambda = \qty{5.32e-7}{m}$}
\find{(a) $\theta$ for $m = 1$ \quad (b) highest $m$ \quad (c) number of spots}
\formula{$d = \dfrac{1}{N}$; \quad $\sin\theta = \dfrac{m\lambda}{d}$ (divide both sides
  of $d\sin\theta = m\lambda$ by $d$); \quad highest $m$ $=$ $\dfrac{d}{\lambda}$ rounded
  down}
\substitute{\eqlines{d &= \frac{1}{\num{5.00e5}\,\unit{m^{-1}}} = \qty{2.00e-6}{m}\\
  \sin\theta &= \frac{1 \times \qty{5.32e-7}{m}}{\qty{2.00e-6}{m}} = 0.266,
  \text{ so } \theta = \ang{15.4}\\
  \frac{d}{\lambda} &= \frac{\qty{2.00e-6}{m}}{\qty{5.32e-7}{m}} = 3.76
  \text{, so the highest order is } m = 3}}
\answer{(a) \final{$\theta = \ang{15.4}$} \quad (b) \final{$m = 3$} (at \ang{52.9}) \quad
  (c) \final{7 spots}: $m = 0$, plus $m = 1, 2, 3$ on each side}
\checkstep{$m = 4$ would need $\sin\theta = 1.064$, which is impossible \cmark.}
\end{workedexample}

\begin{workedexample}{Finding a laser's wavelength from the spots}\label{ch2:we:spots}
A laser shines through a grating with 600 lines per mm onto a wall \qty{1.50}{m} away. The
two first-order spots are \qty{1.21}{m} apart. Find the wavelength of the laser light.
\given{$N = 600 \times 1000 = \num{6.00e5}$ lines per m; \ $L = \qty{1.50}{m}$; \
  first-order spots \qty{1.21}{m} apart, so $m = 1$}
\find{the wavelength $\lambda$}
\formula{$d = \dfrac{1}{N}$; \quad $\tan\theta = \dfrac{y}{L}$; \quad
  $\lambda = \dfrac{d\sin\theta}{m}$ (divide both sides of $d\sin\theta = m\lambda$ by $m$)}
\substitute{\eqlines{d &= \frac{1}{\num{6.00e5}\,\unit{m^{-1}}} = \qty{1.667e-6}{m}\\
  y &= \qty{1.21}{m} \div 2 = \qty{0.605}{m} \quad \text{(half the gap between the spots)}\\
  \tan\theta &= \frac{\qty{0.605}{m}}{\qty{1.50}{m}} = 0.4033\\
  \theta &= \tan^{-1}(0.4033) = \ang{21.97}\\
  \sin\theta &= \sin\ang{21.97} = 0.3741\\
  \lambda &= \frac{\qty{1.667e-6}{m} \times 0.3741}{1} = \qty{6.234e-7}{m}}}
\answer{\final{$\lambda = \qty{6.23e-7}{m} = \qty{623}{nm}$} (3 significant figures)}
\checkstep{\qty{623}{nm} is orange-red light, inside 400--\qty{700}{nm} \cmark. Using
  $\sin\theta \approx \dfrac{y}{L}$ would give \qty{672}{nm}, about 8\,\% too big, because
  \ang{22} is not a small angle.}
\end{workedexample}

\begin{mistakes}
\begin{itemize}
  \item \xmark\ Rounding the highest order up (3.76 $\to$ 4). \cmark\ Round down.
  \item \xmark\ ``Violet is bent most, like in a prism.'' \cmark\ In a grating (and in
        diffraction) red goes to the biggest angle.
  \item \xmark\ Taking $y$ as the distance between the two first-order spots. \cmark\ $y$
        is half of it: from the centre to one spot.
  \item \xmark\ Using $\Delta y = \dfrac{\lambda L}{d}$ for a grating. \cmark\ Grating angles
        are large: use $\tan^{-1}\!\left(\dfrac{y}{L}\right)$, then $d\sin\theta = m\lambda$.
\end{itemize}
\end{mistakes}

\begin{examtip}
\begin{itemize}
  \item ``Explain how a grating makes a spectrum'' (3 marks): cause (each $\lambda$ has
        its own angle) $\to$ link (longer $\lambda$, bigger $\theta$) $\to$ effect (colours
        spread out, violet nearest the white centre).
  \item ``Calculate the wavelength from the spot positions'' (3--4 marks): one mark each
        for $d = \dfrac{1}{N}$ in metres, $\theta$ from $\tan^{-1}\!\left(\dfrac{y}{L}\right)$,
        $d\sin\theta = m\lambda$ rearranged, and the answer with a unit.
\end{itemize}
\end{examtip}

\begin{quickcheck}
  \item A grating has 800 lines per mm. Blue light of \qty{445}{nm} passes through it.
        Find $d$, the first-order angle, and the highest order.
  \item A laser shines through a grating with 400 lines per mm onto a wall \qty{1.20}{m}
        away. The two first-order spots are \qty{0.640}{m} apart. Find the wavelength.
  \item White light passes through a grating. In the first-order spectrum, which colour is
        closest to the central white line? Explain using $d\sin\theta = m\lambda$.
\end{quickcheck}

%% ==================================================================
\section{Polarisation and Malus' law}\label{ch2:sec:polarisation}
\covers{M7.IQ2.04, M7.IQ2.03}

\begin{bigidea}
In light, the electric field vibrates at \ang{90} to the direction of travel: light is a
transverse wave. A polarising filter lets through only one direction of vibration. Only
transverse waves can be polarised, so polarisation tells us what kind of wave light is.
\end{bigidea}

\begin{keywords}
  \kw{Polarised light}{light whose electric field vibrates in one direction only.}
  \kw{Transmission axis}{the direction of vibration that a polarising filter lets through.}
  \kw{Intensity ($I$)}{the power of light falling on each square metre. Unit: \si{W.m^{-2}}.}
\end{keywords}

\begin{refresher}{Transverse and longitudinal waves (Year 11)}
In a \term{transverse wave} the vibration is at \ang{90} to the direction of travel (a rope
shaken up and down). In a \term{longitudinal wave} it is along the direction of travel (a
slinky pushed back and forth, or sound). Light is an electromagnetic wave
(\cref{ch1:chap}): its electric and magnetic fields vibrate at \ang{90} to each other and
to the direction of travel.
\end{refresher}

A rope wave shaken up and down passes through a \emph{vertical} gap in a fence, but a
\emph{horizontal} gap stops it. A back-and-forth (longitudinal) wave would pass both gaps
in the same way. A transverse wave that vibrates in one direction only is
\term{polarised}. This property is called \term{polarisation}. A filter is one way to
polarise a wave. Only transverse waves can be polarised, because only they have a sideways
direction of vibration for a filter to select.

Light from the Sun or a lamp is \term{unpolarised}: its electric field vibrates in every
direction across the beam, changing at random. A \term{polariser} passes only the part of
the vibration along its \term{transmission axis}. The light that comes out is
\term{plane-polarised}: it vibrates in one direction only. A second filter, the
\term{analyser}, passes only the part along its own axis (\cref{ch2:fig:polarisers}).

\begin{figure}
  \centering
  \begin{tikzpicture}[x=0.95cm, y=0.95cm]
    \begin{scope}
      \draw[hscmuted!50, line width=3pt, -{Stealth[length=4mm]}] (-0.9,-1.6) -- (8.9,-1.6);
      \node[text=hscmuted, font=\footnotesize, fill=white, inner sep=1pt] at (3.1,-1.6) {direction of travel};
      \draw[hscmuted, line width=0.5pt] (0,0) circle (0.62);
      \foreach \a in {0,45,90,135} \draw[hscfield, line width=1pt, {Stealth[length=2mm]}-{Stealth[length=2mm]}] (\a:0.5) -- (\a+180:0.5);
      \node[font=\footnotesize, anchor=north] at (0,-0.66) {unpolarised, $I_0$};
      \fill[hscmuted!15] (2.1,0) circle (0.7);
      \draw[hscink, line width=0.6pt] (2.1,0) circle (0.7);
      \draw[hscink, line width=1.4pt] (2.1,-0.7) -- (2.1,0.7);
      \node[font=\scriptsize, anchor=west] at (2.15,0.22) {axis};
      \node[font=\footnotesize, anchor=south] at (2.1,0.72) {polariser};
      \draw[hscmuted, line width=0.5pt] (4.1,0) circle (0.62);
      \draw[hscfield, line width=1.2pt, {Stealth[length=2mm]}-{Stealth[length=2mm]}] (4.1,-0.5) -- (4.1,0.5);
      \node[font=\footnotesize, anchor=north] at (4.1,-0.66) {$\tfrac{1}{2}I_0$};
      \fill[hscmuted!15] (6.1,0) circle (0.7);
      \draw[hscink, line width=0.6pt] (6.1,0) circle (0.7);
      \draw[guide] (6.1,-0.7) -- (6.1,0.7);
      \draw[hscink, line width=1.4pt] ($(6.1,0)+(-55:0.7)$) -- ($(6.1,0)+(125:0.7)$);
      \draw[hscink, line width=0.5pt] (6.1,0.45) arc (90:125:0.45);
      \node[font=\footnotesize] at ($(6.1,0)+(105:0.62)$) {$\theta$};
      \node[font=\footnotesize, anchor=south] at (6.1,0.72) {analyser};
      \draw[hscmuted, line width=0.5pt] (8.1,0) circle (0.62);
      \draw[hscfield, line width=1.2pt, {Stealth[length=2mm]}-{Stealth[length=2mm]}] ($(8.1,0)+(-55:0.41)$) -- ($(8.1,0)+(125:0.41)$);
      \node[font=\footnotesize, anchor=north] at (8.1,-0.66) {$\tfrac{1}{2}I_0\cos^2\theta$};
      \node[font=\footnotesize\bfseries] at (4.0,-2.2) {(a) looking along the beam at each stage};
    \end{scope}
    \begin{scope}[xshift=11.3cm, yshift=-0.55cm]
      \draw[guide] (0,-0.6) -- (0,1.75);
      \draw[hscink, line width=0.8pt] (-55:0.7) -- (125:1.85) node[left, font=\footnotesize] {analyser axis};
      \draw[hscfield, line width=1.2pt, ->] (0,0) -- (0,1.5) node[right, font=\footnotesize, text=hscfield] {$A$};
      \draw[hscmuted, dashed, line width=0.5pt] (0,1.5) -- (125:1.229);
      \draw[hscforce, line width=1.6pt, ->] (0,0) -- (125:1.229);
      \draw[hscink, line width=0.5pt] (0,0.55) arc (90:125:0.55);
      \node[font=\footnotesize] at (108:0.75) {$\theta$};
      \node[text=hscforce, font=\footnotesize, anchor=north east] at (125:0.75) {$A\cos\theta$};
      \node[font=\footnotesize\bfseries, align=center] at (0,-1.65) {(b) the part that passes};
    \end{scope}
  \end{tikzpicture}
  \caption{(a) Each circle shows the light as if the beam were coming straight towards
    you. The dark line on each filter is its transmission axis. The polariser passes only
    the vertical vibration, so $I_\mathrm{max} = \tfrac{1}{2}I_0$ reaches the analyser. The
    analyser, at $\theta$ to the polariser, passes only the part along its own axis.
    (b) A vibration of amplitude $A$ has a part $A\cos\theta$ along the analyser's axis.}
  \label{ch2:fig:polarisers}
\end{figure}

\begin{formulacard}[notonsheet]{Unpolarised light through one polariser}
  \cardformula{I = \tfrac{1}{2}I_0}
  \begin{symbols}
    \sym{$I_0$}{intensity of the unpolarised light arriving}{\si{W.m^{-2}}}
    \sym{$I$}{intensity coming out of the filter}{\si{W.m^{-2}}}
  \end{symbols}
  \inwords{an ideal polariser passes half of unpolarised light, at any angle.}
  \whentouse{unpolarised light (sunlight, a lamp) meets the first polariser.}
  \whennottouse{the light is already polarised: use Malus' law.}
  \rearranged{$I_0 = 2I$}
  \tensecond{\qty{50}{W.m^{-2}} of sunlight through one polariser gives \qty{25}{W.m^{-2}}.}
  \derive{split every vibration into two parts: one along the transmission axis and one at
    \ang{90} to it. In unpolarised light no direction is special. So, on average, the two
    parts carry the same energy. The filter passes the part along its axis and blocks the
    other part. So half of the intensity passes.}
  \trap{do not use $\cos^2\theta$ at the first filter for unpolarised light. Its angle does
    not matter.}
\end{formulacard}

\begin{formulacard}{Malus' law}
  \cardformula{I = I_\mathrm{max}\cos^2\theta}
  \begin{symbols}
    \sym{$I$}{intensity that passes through the filter}{\si{W.m^{-2}}}
    \sym{$I_\mathrm{max}$}{intensity of the polarised light arriving (what passes at $\theta = 0$)}{\si{W.m^{-2}}}
    \sym{$\theta$}{angle between the light's vibration and the filter's transmission axis}{degree (\si{\degree})}
  \end{symbols}
  \inwords{the further the filter is turned from the light's vibration, the less passes. At
    \ang{90}, none passes.}
  \whentouse{light that is \emph{already} plane-polarised meets a polariser.}
  \whennottouse{unpolarised light: use $I = \frac{1}{2}I_0$.}
  \rearranged{$I_\mathrm{max} = \dfrac{I}{\cos^2\theta}$ \qquad
    $\cos\theta = \sqrt{\dfrac{I}{I_\mathrm{max}}}$ \qquad
    $\theta = \cos^{-1}\!\sqrt{\dfrac{I}{I_\mathrm{max}}}$}
  \tensecond{$I_\mathrm{max} = \qty{40}{W.m^{-2}}$, $\theta = \ang{60}$:
    $I = 40 \times (\cos\ang{60})^2 = 40 \times 0.25 = \qty{10}{W.m^{-2}}$.}
  \derive{only the part $A\cos\theta$ of the amplitude $A$ (the size of the vibration) passes
    (\cref{ch2:fig:polarisers}b). Intensity is proportional to amplitude squared (double the
    amplitude gives four times the intensity). So
    $\dfrac{I}{I_\mathrm{max}} = \dfrac{(A\cos\theta)^2}{A^2} = \cos^2\theta$.}
  \trap{$\cos^2\theta$ means $(\cos\theta)^2$. Measure $\theta$ from the \emph{previous}
    filter, not from the vertical.}
\end{formulacard}

\begin{toolkit}{Gradient of a straight-line graph}
Gradient $=$ rise $\div$ run $= \dfrac{\text{change in the vertical value}}{\text{change in
the horizontal value}}$. Use two points far apart on the line of best fit, not two data
points. Its unit is the vertical unit $\div$ the horizontal unit. A straight line through
the origin means the vertical quantity is proportional to the horizontal one.
\end{toolkit}

\Cref{ch2:fig:malus} shows the two graphs you must know. The graph of $I$ against $\theta$
is a curve: greatest at \ang{0} and \ang{180}, zero at \ang{90}. The graph of $I$ against
$\cos^2\theta$ is a straight line through the origin with gradient $I_\mathrm{max}$. The
points in (b) are the sample results of the Malus' law investigation
(\cref{ch2:sec:experiments}).

\begin{figure}
  \centering
  \begin{tikzpicture}
    \begin{axis}[hsc, width=0.47\linewidth, height=0.31\linewidth,
        xmin=0, xmax=180, ymin=0, ymax=1.15, xtick={0,45,90,135,180},
        ytick={0,0.5,1}, yticklabels={0, $\tfrac{1}{2}I_\mathrm{max}$, $I_\mathrm{max}$},
        xlabel={$\theta$ (\si{\degree})}, ylabel={$I$}, title={(a) theory: $I$ against $\theta$},
        title style={font=\footnotesize\bfseries}]
      \addplot[hscseriesA, line width=1.2pt, domain=0:180, samples=91] {cos(x)^2};
    \end{axis}
  \end{tikzpicture}\hfill
  \begin{tikzpicture}
    \begin{axis}[hsc, width=0.47\linewidth, height=0.31\linewidth,
        xmin=0, xmax=1.05, ymin=0, ymax=13.5, xtick={0,0.25,0.5,0.75,1},
        xlabel={$\cos^2\theta$}, ylabel={$I$ (\si{W.m^{-2}})},
        title={(b) data: $I$ against $\cos^2\theta$}, title style={font=\footnotesize\bfseries},
        legend style={at={(0.98,0.03)}, anchor=south east, font=\scriptsize}]
      \addplot[hscseriesA, only marks, mark=*, mark size=1.8pt]
        coordinates {(1.000,12.0) (0.933,11.3) (0.750,8.9) (0.500,6.1) (0.250,3.1) (0.067,0.9) (0.000,0.1)};
      \addlegendentry{\ang{0} to \ang{90}}
      \addplot[hscseriesC, only marks, mark=o, mark size=2.6pt, line width=0.8pt]
        coordinates {(0.067,0.9) (0.250,3.0) (0.500,6.0) (0.750,9.0) (0.933,11.2) (1.000,12.0)};
      \addlegendentry{\ang{105} to \ang{180}}
      \addplot[hscseriesB, dashed, line width=1pt, domain=0:1.05, samples=2] {12.02*x};
      \node[anchor=west, font=\footnotesize, text=hscink] at (axis cs:0.02,11.8)
        {gradient $= I_\mathrm{max} = \qty{12.0}{W.m^{-2}}$};
    \end{axis}
  \end{tikzpicture}
  \caption{Malus' law. (a) $I$ against $\theta$ is a $\cos^2$ curve. (b) $I$ against
    $\cos^2\theta$ is a straight line through the origin.}
  \label{ch2:fig:malus}
\end{figure}

\paragraph{Three filters.} Two filters at \ang{90} (``crossed'') block all the light. Slide
a third filter at \ang{45} \emph{between} them, and light gets through! Each filter
re-polarises the light along its own axis. So the vibration is turned step by step
(\cref{ch2:we:three}).

\paragraph{Uses.} Glare is light reflected from water or a road. It is partly polarised
horizontally. Polarising sunglasses have a vertical transmission axis, so they block much
of the glare. LCD screens and 3D-cinema glasses also use polarisers.

\paragraph{Why polarisation matters for the model of light.}
\begin{itemize}
  \item Only transverse waves can be polarised, so light must be a \emph{transverse} wave.
  \item Huygens' wave was longitudinal. A longitudinal wave cannot be polarised, so his
        model had to change. Young (1817) and Fresnel (1821) proposed that light waves are
        transverse.
  \item Newton's particle model had no natural explanation. (Newton could only guess that
        light rays have ``sides''.) A transverse wave explains polarisation, and Malus' law
        follows from the part $A\cos\theta$ of the vibration.
  \item Maxwell's theory (1865) predicts that light is an electromagnetic wave with fields at
        \ang{90} to the travel: a transverse wave. Polarisation fits it exactly.
\end{itemize}
So polarisation was highly significant. With interference, it completed the classical
model: light is a transverse electromagnetic wave. (\Cref{ch3:chap} shows that even this
model is not the whole story.)

\begin{workedexample}{Sunlight through two filters}\label{ch2:we:two}
Unpolarised sunlight of \qty{600}{W.m^{-2}} passes through a polariser and then an
analyser at \ang{25.0} to it. Find the intensity after each filter.
\given{$I_0 = \qty{600}{W.m^{-2}}$ (unpolarised), \ $\theta = \ang{25.0}$ between the filters}
\find{$I_1$ after the polariser, \ $I_2$ after the analyser}
\formula{$I_1 = \tfrac{1}{2}I_0$ \quad then \quad $I_2 = I_\mathrm{max}\cos^2\theta$ with
  $I_\mathrm{max} = I_1$}
\substitute{\eqlines{I_1 &= \tfrac{1}{2} \times \qty{600}{W.m^{-2}} = \qty{300}{W.m^{-2}}\\
  I_2 &= \qty{300}{W.m^{-2}} \times (\cos\ang{25.0})^2 = 300 \times 0.8214 = \qty{246.4}{W.m^{-2}}}}
\answer{\final{$I_1 = \qty{300}{W.m^{-2}}$} and \final{$I_2 = \qty{246}{W.m^{-2}}$}}
\checkstep{\ang{25.0} is much less than \ang{90}, so most of the light
  ($\cos^2\ang{25.0} = 0.82$, i.e.\ 82\,\%) passes the analyser. $I_2$ is less than $I_1$,
  as it must be \cmark.}
\end{workedexample}

\begin{workedexample}{Three filters let more light through than two}\label{ch2:we:three}
Unpolarised light of \qty{480}{W.m^{-2}} meets a vertical polariser and then a horizontal
one. (a) Find the intensity that comes out. (b) A third filter at \ang{30} to the vertical
is placed between them. Find the new intensity. (c) Which angle of the middle filter
passes the most light?
\given{$I_0 = \qty{480}{W.m^{-2}}$. Filters at \ang{0} and \ang{90} to the vertical. In (b)
  a middle filter is at \ang{30}, so the angle between it and the last filter is
  $\ang{90} - \ang{30} = \ang{60}$.}
\find{(a) $I$ with two filters \quad (b) $I$ with three \quad (c) the best middle angle:
  try several angles (\ang{30} is done in (b))}
\formula{first filter $I = \tfrac{1}{2}I_0$; each later filter $I = I_\mathrm{max}\cos^2\theta$,
  with $\theta$ from the filter before it}
\substitute{\eqlines{\text{(a) } I &= \tfrac{1}{2}(480) \times (\cos\ang{90})^2
  = 240 \times 0 = \qty{0}{W.m^{-2}}\\
  \text{(b) } I_1 &= \tfrac{1}{2} \times 480 = \qty{240}{W.m^{-2}}\\
  I_2 &= 240 \times (\cos\ang{30})^2 = 240 \times 0.75 = \qty{180}{W.m^{-2}}\\
  I_3 &= 180 \times (\cos\ang{60})^2 = 180 \times 0.25 = \qty{45}{W.m^{-2}}}\par
  (c) Middle filter at \ang{60}, so \ang{30} from the last filter:\par
  \eqlines{I_3 &= 240 \times (\cos\ang{60})^2 \times (\cos\ang{30})^2\\
  &= 240 \times 0.25 \times 0.75 = \qty{45}{W.m^{-2}}}\par
  Middle filter at \ang{45}, so \ang{45} from the last filter:\par
  \eqlines{I_3 &= 240 \times (\cos\ang{45})^2 \times (\cos\ang{45})^2\\
  &= 240 \times 0.500 \times 0.500 = \qty{60}{W.m^{-2}}}\par
  In the same way, \ang{15} and \ang{75} each give \qty{15}{W.m^{-2}}.}
\answer{(a) \final{\qty{0}{W.m^{-2}}} \quad (b) \final{\qty{45}{W.m^{-2}}} \quad
  (c) \final{\ang{45}} (\qty{60}{W.m^{-2}}, the most)}
\checkstep{\ang{30} and \ang{60} give the same \qty{45}{W.m^{-2}}. Half-way, at \ang{45},
  gives more. Two equal \ang{45} steps turn the vibration through \ang{90} with the least
  loss \cmark.}
\end{workedexample}

\begin{mistakes}
\begin{itemize}
  \item \xmark\ $\cos^2\theta$ for unpolarised light. \cmark\ The first filter passes half:
        $I = \frac{1}{2}I_0$.
  \item \xmark\ $\cos\theta$ instead of $\cos^2\theta$. \cmark\ Square the cosine.
  \item \xmark\ Measuring every angle from the vertical. \cmark\ Use the angle from the
        filter before.
  \item \xmark\ ``Polarisation proves light is electromagnetic.'' \cmark\ It shows light is
        \emph{transverse}; Maxwell's theory explains why.
\end{itemize}
\end{mistakes}

\begin{examtip}
\begin{itemize}
  \item Multiple choice ``Which graph shows Malus' law?'': $I$ against $\cos^2\theta$ is a
        straight line through the origin (\cref{ch2:fig:malus}b).
  \item ``Evaluate the significance of polarisation'' (4--5 marks): make the four points
        in the list above, then finish with a judgement (``highly significant, because
        \ldots'').
  \item With several filters, show the intensity after \emph{each} filter.
\end{itemize}
\end{examtip}

\begin{quickcheck}
  \item Polarised light of \qty{20.0}{W.m^{-2}} meets a filter at \ang{60.0} to its
        vibration. Calculate the intensity that passes.
  \item A lamp's light passes one polarising filter with its axis at \ang{30} to the
        vertical. \qty{12}{W.m^{-2}} comes out. Find the lamp's intensity before the filter.
  \item Polarised light meets a polarising filter. Only 40\,\% of it passes
        ($I = 0.40\,I_\mathrm{max}$). Find the angle between the light's vibration and the
        filter's axis.
  \item Unpolarised light of \qty{200}{W.m^{-2}} passes through three filters with axes at
        \ang{0}, \ang{20} and \ang{90} to the vertical. Find the intensity that comes out.
  \item Explain why sound waves cannot be polarised, but light can.
\end{quickcheck}

%% ==================================================================
\section{Key experiments}\label{ch2:sec:experiments}
\covers{M7.IQ2.01, M7.IQ2.02, M7.IQ2.03, M7.IQ2.04}

This section describes the historical experiments and the three required investigations
for this chapter.

\begin{experiment}{Foucault: is light faster or slower in water? (1850)}
\aim{To compare the speed of light in water and in air, and so test Newton's and Huygens'
  predictions for refraction.}
\apparatus{a light source; a glass plate (it lets light through, and also reflects some of
  the returning light into the eyepiece); a fast-spinning mirror; a distant fixed mirror; a
  tube of water; and an eyepiece with a scale (\cref{ch2:fig:foucault}).
  \begin{figure}
    \centering
    \begin{tikzpicture}[x=0.72cm, y=0.72cm]
      \fill[hsctip!70] (2.4,-3.0) circle (0.2); \draw[hscink, line width=0.5pt] (2.4,-3.0) circle (0.2);
      \node[font=\footnotesize, anchor=west] at (2.65,-3.0) {source};
      \draw[ray] (2.4,-2.8) -- (2.4,0);
      \draw[hscvel!70, line width=1.6pt] (2.12,-1.88) -- (2.68,-1.32);
      \node[font=\footnotesize, anchor=west] at (2.75,-1.6) {glass plate};
      \fill[hscvel!15] (4.6,-0.28) rectangle (8.2,0.28);
      \draw[hscvel, line width=0.6pt] (4.6,-0.28) rectangle (8.2,0.28);
      \node[font=\footnotesize, anchor=north] at (6.4,-0.35) {tube of water (second trial)};
      \draw[ray] (2.4,0) -- (10.0,0);
      \draw[hscink, line width=2.2pt] (2.12,-0.28) -- (2.68,0.28);
      \fill[hscink] (2.4,0) circle (1.6pt);
      \draw[->, hscink, line width=0.6pt] (2.95,0.25) arc (20:110:0.6);
      \node[font=\footnotesize, anchor=south] at (2.4,0.75) {spinning mirror};
      \draw[hscink, line width=2.2pt] ($(2.4,0)+(-5:7.6)$) arc[start angle=-5, end angle=5, radius=7.6];
      \node[font=\footnotesize, anchor=west] at (10.15,0) {fixed mirror};
      \draw[hscray, line width=0.7pt, dashed, ->] (2.4,-1.6) -- (0.75,-1.72);
      \draw[hscray, line width=0.7pt, dashed, ->] (2.4,-1.6) -- (0.75,-2.05);
      \draw[hscink, line width=0.6pt] (0.65,-2.4) -- (0.65,-1.2);
      \foreach \y in {-2.3,-2.2,...,-1.25} \draw[hscink, line width=0.4pt] (0.65,\y) -- (0.55,\y);
      \node[font=\scriptsize, anchor=east] at (0.5,-1.72) {air};
      \node[font=\scriptsize, anchor=east] at (0.5,-2.05) {water};
      \node[font=\footnotesize, anchor=north, align=center] at (0.45,-2.45) {scale in\\the eyepiece};
    \end{tikzpicture}
    \caption{Foucault's method (simplified). The spinning mirror reflects the light to the
      fixed mirror, and the light comes back the same way. A bigger shift on the scale
      means a longer travel time.}
    \label{ch2:fig:foucault}
  \end{figure}}
\method{\begin{enumerate}
  \item Light reflects off the spinning mirror to the fixed mirror and back. The mirror
        turns a little during the round trip, so the returning image is shifted along the
        scale. Measure the shift.
  \item Put the water tube in the path, keep the spin rate the same, and measure the shift
        again.
\end{enumerate}}
\results{the shift was bigger with water in the path.}
\conclusion{the light took longer, so light is slower in water than in air.}
\significance{Newton's model needed light to be faster in water, so his explanation of
  refraction was wrong. Huygens' model agreed. The result disproved one model; it did not
  prove the other on its own.}
\validity{\begin{itemize}
  \item \textbf{Validity:} the same apparatus and spin rate in both trials, so only the
        material in the path changes.
  \item \textbf{Reliability:} repeat each measurement and average.
  \item \textbf{Accuracy:} the shift is tiny, so the mirror must spin very fast at a steady,
        known rate.
\end{itemize}}
\end{experiment}

\begin{experiment}{Diffraction and interference with lasers}
\aim{(A) To see how diffraction depends on the slit width and the wavelength
  (qualitative). (B) To measure the wavelength of a red laser with a double slit and a
  grating, using $d\sin\theta = m\lambda$ (quantitative). Young first did the double-slit
  part of (B), with sunlight, in about 1801.}
\apparatus{a red laser (label: $\lambda = \qty{650}{nm}$) and a green laser (label:
  $\lambda = \qty{532}{nm}$), both low power; single slits \qty{0.080}{mm},
  \qty{0.040}{mm} and \qty{0.020}{mm} wide; a double slit ($d = \qty{0.250}{mm}$); a
  grating (300 lines per mm); a screen; a tape measure; a darkened room
  (\cref{ch2:fig:diffsetup}). Young used sunlight through a pinhole and a card with two
  slits (\cref{ch2:fig:young}).
  \begin{figure}
    \centering
    \begin{tikzpicture}[x=0.85cm, y=0.7cm]
      \fill[hscmuted!35] (0,-0.25) rectangle (1.4,0.25);
      \draw[hscink, line width=0.6pt] (0,-0.25) rectangle (1.4,0.25);
      \node[font=\footnotesize] at (0.7,0) {laser};
      \draw[hscforce, line width=1.2pt] (1.4,0) -- (2.6,0);
      \fill[hscink!70] (2.6,-0.9) rectangle (2.7,0.9);
      \node[font=\footnotesize, anchor=south] at (2.65,0.95) {slit, double slit or grating};
      \fill[hscforce!25] (2.7,0) -- (8.0,1.35) -- (8.0,-1.35) -- cycle;
      \draw[hscforce, line width=0.8pt] (2.7,0) -- (8.0,0);
      \draw[hscforce, line width=0.6pt, dashed] (2.7,0) -- (8.0,0.85);
      \fill[hscink!88] (8.0,-1.5) rectangle (8.25,1.5);
      \foreach \yy/\h in {0/0.2, 0.85/0.1, -0.85/0.1} \fill[hscforce!85!white] (8.0,\yy-\h) rectangle (8.25,\yy+\h);
      \node[font=\footnotesize, anchor=south] at (8.12,1.55) {screen};
      \draw[|<->|, line width=0.5pt] (2.65,-1.25) -- (8.0,-1.25) node[midway, below, font=\footnotesize] {$L$};
      \draw[|<->|, line width=0.5pt] (8.55,0) -- (8.55,0.85) node[midway, right, font=\footnotesize] {$y$};
    \end{tikzpicture}
    \caption{Set-up for the laser investigations (seen from above).}
    \label{ch2:fig:diffsetup}
  \end{figure}}
\method{\begin{enumerate}
  \item (A) Shine each laser through each single slit, with $L = \qty{2.00}{m}$. Photograph
        each pattern next to a ruler and measure the central band.
  \item (B) Double slit: measure across 10 fringe spacings and divide by 10 to get
        $\Delta y$. Do this for $L$ from 1.00 to \qty{3.00}{m}, and plot $\Delta y$ against $L$.
  \item (B) Grating, $L = \qty{1.000}{m}$: measure the distance between the two first-order
        spots and halve it to get $y$.
  \item Safety: never look into the beam or its reflections; keep it below eye level.
\end{enumerate}}
\results{(A) Width of the central band (sample results):
  \begin{center}\small
  \begin{tabular}{@{}lcc@{}}
    \toprule
    Slit width (\si{mm}) & Red (\si{mm}) & Green (\si{mm})\\
    \midrule
    0.080 & 33 & 27\\
    0.040 & 65 & 53\\
    0.020 & 130 & 106\\
    \bottomrule
  \end{tabular}
  \end{center}
  Every pattern looked like \cref{ch2:fig:singleslit}.\par
  (B) Double slit (sample results):
  \begin{center}\small
  \begin{tabular}{@{}lccccc@{}}
    \toprule
    $L$ (\si{m}) & 1.00 & 1.50 & 2.00 & 2.50 & 3.00\\
    $\Delta y$ (\si{mm}) & 2.56 & 3.94 & 5.15 & 6.58 & 7.72\\
    \bottomrule
  \end{tabular}
  \end{center}
  \begin{figure}
    \centering
    \begin{tikzpicture}
      \begin{axis}[hsc, width=0.56\linewidth, height=0.26\linewidth,
          xmin=0, xmax=3.3, ymin=0, ymax=9, xtick={0,0.5,...,3},
          xticklabel style={/pgf/number format/.cd, fixed, precision=1},
          xlabel={$L$ (\si{m})}, ylabel={$\Delta y$ (\si{mm})}]
        \addplot[hscseriesA, only marks, mark=*, mark size=1.8pt]
          coordinates {(1.00,2.56) (1.50,3.94) (2.00,5.15) (2.50,6.58) (3.00,7.72)};
        \addplot[hscseriesB, dashed, line width=1pt, domain=0:3.3, samples=2] {2.5947*x};
        \node[anchor=west, font=\footnotesize, text=hscink] at (axis cs:1.55,1.6)
          {gradient $= \qty{2.59}{mm.m^{-1}}$};
      \end{axis}
    \end{tikzpicture}
    \caption{Fringe spacing against screen distance (double slit).}
    \label{ch2:fig:dygraph}
  \end{figure}
  Analysis, step by step:\par\nopagebreak
  \begin{enumerate}[beginpenalty=10000]
    \item Double slit: the line of best fit goes through $(0, 0)$ and
          $(\qty{3.00}{m}, \qty{7.78}{mm})$. Gradient $=$ rise $\div$ run
          $= \dfrac{\qty{7.78}{mm} - 0}{\qty{3.00}{m} - 0} = \qty{2.59}{mm.m^{-1}}$.
          Change mm to m ($\times 10^{-3}$): \qty{2.59e-3}{m} per m. The units cancel, so
          the gradient is \num{2.59e-3} (no unit).
    \item $\Delta y = \dfrac{\lambda}{d}L$, so the gradient is $\dfrac{\lambda}{d}$.
          Multiply both sides by $d$: $\lambda = \text{gradient} \times d =
          \num{2.59e-3} \times \qty{2.50e-4}{m} = \qty{6.48e-7}{m}$.
    \item Grating: $d = \dfrac{1}{N} = \dfrac{1}{300 \times \qty{1000}{m^{-1}}}
          = \dfrac{1}{\num{3.00e5}\,\unit{m^{-1}}} = \qty{3.333e-6}{m}$.
    \item The two first-order spots were \qty{397}{mm} apart, so
          $y = \qty{397}{mm} \div 2 = \qty{198.5}{mm} = \qty{0.1985}{m}$.
          $\tan\theta = \dfrac{y}{L} = \dfrac{\qty{0.1985}{m}}{\qty{1.000}{m}} = 0.1985$,
          so $\theta = \tan^{-1}(0.1985) = \ang{11.23}$.
    \item First order, $m = 1$: $\sin\ang{11.23} = 0.1947$, so
          $\lambda = \dfrac{d\sin\theta}{m} = \dfrac{\qty{3.333e-6}{m} \times 0.1947}{1}
          = \qty{6.49e-7}{m}$.
  \end{enumerate}}
\conclusion{(A) Light spreads more through a narrower slit and with a longer wavelength,
  as waves do. (B) Light from two slits interferes: bright where $\delta = m\lambda$, dark
  where $\delta = \left(m - \frac{1}{2}\right)\lambda$. The double slit gives
  $\lambda = \qty{648}{nm}$ and the grating gives \qty{649}{nm}. Both are within 0.3\,\% of
  the label value (\qty{650}{nm}).}
\significance{(A) Light spreads through narrow slits, just as water and sound waves do. It
  spreads noticeably only when the slit is very narrow, which is why everyday shadows look
  sharp. (B) Particles cannot cancel to make dark fringes, so Young's result was strong
  evidence for waves. It also gave the first wavelengths of colours.}
\validity{\begin{itemize}
  \item \textbf{Validity:} change one variable at a time. A laser (or Young's single
        source) makes the slits coherent. Use $\Delta y = \dfrac{\lambda L}{d}$ only at small
        angles (here less than \ang{2}); for the grating (\ang{11}) use $\tan^{-1}$, then $\sin$.
  \item \textbf{Reliability:} repeat readings; measure across 10 spacings, on both sides of
        the centre; use a line through 5 points.
  \item \textbf{Accuracy:} percentage error $= \dfrac{648 - 650}{650} \times 100\,\% =
        -0.31\,\%$ (double slit) and $\dfrac{649 - 650}{650} \times 100\,\% = -0.15\,\%$
        (grating). Sharp grating spots and a bigger $L$ make the percentage
        \emph{uncertainty} in $y$ smaller, which improves accuracy.
\end{itemize}}
\end{experiment}

\begin{experiment}{Testing Malus' law}
\aim{To test Malus' law, $I = I_\mathrm{max}\cos^2\theta$, for plane-polarised light.}
\apparatus{a lamp; a polariser and an analyser with an angle scale; a light sensor (a lux
  meter or a data logger; its reading is proportional to intensity, so the graphs have the
  same shape); a darkened room (\cref{ch2:fig:malussetup}).
  \begin{figure}
    \centering
    \begin{tikzpicture}[x=0.85cm, y=0.6cm]
      \fill[hsctip!70] (0,0) circle (0.25); \draw[hscink, line width=0.5pt] (0,0) circle (0.25);
      \foreach \a in {0,45,...,315} \draw[hsctip, line width=0.6pt] (\a:0.32) -- (\a:0.46);
      \node[font=\footnotesize, anchor=north] at (0,-0.55) {lamp};
      \draw[hscforce!60, line width=3pt] (0.5,0) -- (6.6,0);
      \fill[hscmuted!20] (2.2,0) ellipse (0.18 and 0.8); \draw[hscink, line width=0.6pt] (2.2,0) ellipse (0.18 and 0.8);
      \node[font=\footnotesize, anchor=north] at (2.2,-0.85) {polariser (fixed)};
      \fill[hscmuted!20] (4.6,0) ellipse (0.18 and 0.8); \draw[hscink, line width=0.6pt] (4.6,0) ellipse (0.18 and 0.8);
      \draw[->, hscink, line width=0.6pt] (4.2,0.95) arc (150:30:0.45);
      \node[font=\footnotesize, anchor=south] at (4.6,1.2) {turn by $\theta$};
      \node[font=\footnotesize, anchor=north] at (4.6,-0.85) {analyser};
      \fill[hscmuted!40] (6.6,-0.35) rectangle (7.3,0.35); \draw[hscink, line width=0.6pt] (6.6,-0.35) rectangle (7.3,0.35);
      \node[font=\footnotesize, anchor=north] at (6.95,-0.45) {sensor};
      \draw[hscink, line width=0.6pt] (7.3,0) -- (7.9,0);
      \draw[hscink, line width=0.6pt, fill=white] (7.9,-0.4) rectangle (10.1,0.4);
      \node[font=\footnotesize] at (9.0,0) {data logger};
    \end{tikzpicture}
    \caption{Set-up for testing Malus' law.}
    \label{ch2:fig:malussetup}
  \end{figure}}
\method{\begin{enumerate}
  \item Line up the lamp, filters and sensor in a dark room. With the lamp off, record the
        background reading and subtract it from every later reading.
  \item Turn the analyser until the reading is greatest: this is $\theta = \ang{0}$. Then
        turn it in \ang{15} steps from \ang{0} to \ang{180}, recording $I$.
  \item Plot $I$ against $\theta$, and $I$ against $\cos^2\theta$.
\end{enumerate}}
\results{\begin{center}\small
  \begin{tabular}{@{}lccccccc@{}}
    \toprule
    $\theta$ (\si{\degree}) & 0 & 15 & 30 & 45 & 60 & 75 & 90\\
    $\cos^2\theta$ & 1.000 & 0.933 & 0.750 & 0.500 & 0.250 & 0.067 & 0.000\\
    $I$ (\si{W.m^{-2}}) & 12.0 & 11.3 & 8.9 & 6.1 & 3.1 & 0.9 & 0.1\\
    \midrule
    $\theta$ (\si{\degree}) & 105 & 120 & 135 & 150 & 165 & 180 & \\
    $\cos^2\theta$ & 0.067 & 0.250 & 0.500 & 0.750 & 0.933 & 1.000 & \\
    $I$ (\si{W.m^{-2}}) & 0.9 & 3.0 & 6.0 & 9.0 & 11.2 & 12.0 & \\
    \bottomrule
  \end{tabular}
  \end{center}
  These are sample results (background subtracted). The graph of $I$ against $\theta$ has
  the shape of \cref{ch2:fig:malus}a: the readings at $\theta$ and $\ang{180} - \theta$
  match. The graph of $I$ against $\cos^2\theta$ is \cref{ch2:fig:malus}b.}
\conclusion{the graph of $I$ against $\cos^2\theta$ is a straight line through the origin.
  Two points on the line of best fit are $(0, 0)$ and $(1.00, \qty{12.0}{W.m^{-2}})$, so
  the gradient is $\dfrac{12.0 - 0}{1.00 - 0} = \qty{12.0}{W.m^{-2}} = I_\mathrm{max}$. So
  $I \propto \cos^2\theta$ and Malus' law is supported. Almost no light passes at
  \ang{90}.}
\significance{light can be polarised, so it is a transverse wave, as Maxwell's theory predicts.}
\validity{\begin{itemize}
  \item \textbf{Validity:} change only $\theta$; keep the lamp, distances and sensor fixed;
        subtract background light.
  \item \textbf{Reliability:} repeat the full turn and average. The readings at $\theta$
        and $\ang{180} - \theta$ should match.
  \item \textbf{Accuracy:} read the angle scale carefully, as $I$ changes fastest near
        \ang{45}. Real filters leak a little, so the \ang{90} reading is just above zero.
\end{itemize}}
\end{experiment}
